% Learn fprettier: A Guide for Developers
%
% Build with:   tectonic learn-fprettier.tex
% or:           latexmk -pdf learn-fprettier.tex
%
% The book describes fprettier 0.5.0 (plus the unreleased fixes on main as of
% October 2026). Every before/after example was produced by running that
% build of fprettier on the input shown.

\documentclass[11pt,a4paper,openany]{book}

% ---------------------------------------------------------------------------
% Engine-dependent fonts: tectonic is XeTeX-based, latexmk -pdf is pdfTeX
% ---------------------------------------------------------------------------
\usepackage{iftex}
\ifPDFTeX
  \usepackage[T1]{fontenc}
  \usepackage[utf8]{inputenc}
  \usepackage{lmodern}
\else
  \usepackage{fontspec}
\fi

\usepackage[margin=2.6cm,headheight=14pt]{geometry}
\usepackage{microtype}
\usepackage{xcolor}
\usepackage{amssymb}
\usepackage{graphicx}
\usepackage{booktabs}
\usepackage{tabularx}
\usepackage{array}
\usepackage{enumitem}
\usepackage{fancyhdr}
\usepackage{titlesec}
\usepackage{upquote}
\usepackage{listings}
\usepackage[most]{tcolorbox}
\usepackage{tikz}
\usetikzlibrary{arrows.meta,positioning,shapes.geometric,fit,calc,backgrounds}
\usepackage[hidelinks]{hyperref}

% ---------------------------------------------------------------------------
% Colours
% ---------------------------------------------------------------------------
\definecolor{ink}{HTML}{1F2933}
\definecolor{accent}{HTML}{2B6CB0}
\definecolor{accentlight}{HTML}{E6F0FA}
\definecolor{codebg}{HTML}{F6F8FA}
\definecolor{codeframe}{HTML}{D0D7DE}
\definecolor{keywordc}{HTML}{8250DF}
\definecolor{commentc}{HTML}{6E7781}
\definecolor{stringc}{HTML}{0A3069}
\definecolor{beforec}{HTML}{FFF5F5}
\definecolor{afterc}{HTML}{F0FFF4}
\definecolor{exframe}{HTML}{B7791F}
\definecolor{exbg}{HTML}{FFFBEB}
\definecolor{warnframe}{HTML}{C53030}
\definecolor{warnbg}{HTML}{FFF5F5}

% ---------------------------------------------------------------------------
% Headings and page style
% ---------------------------------------------------------------------------
\titleformat{\chapter}[display]
  {\normalfont\sffamily\huge\bfseries\color{ink}}
  {\color{accent}\large\chaptertitlename\ \thechapter}{8pt}{\Huge}
\titlespacing*{\chapter}{0pt}{20pt}{24pt}
\titleformat{\section}{\normalfont\sffamily\Large\bfseries\color{ink}}{\color{accent}\thesection}{1em}{}
\titleformat{\subsection}{\normalfont\sffamily\large\bfseries\color{ink}}{\color{accent}\thesubsection}{1em}{}
\titleformat{\subsubsection}{\normalfont\sffamily\normalsize\bfseries\color{ink}}{}{0em}{}

\pagestyle{fancy}
\fancyhf{}
\fancyhead[LE,RO]{\sffamily\small\thepage}
\fancyhead[RE]{\sffamily\small\nouppercase{\leftmark}}
\fancyhead[LO]{\sffamily\small\nouppercase{\rightmark}}
\renewcommand{\headrulewidth}{0.4pt}
\fancypagestyle{plain}{\fancyhf{}\fancyfoot[C]{\sffamily\small\thepage}\renewcommand{\headrulewidth}{0pt}}

\setlength{\parskip}{0.45em}
\setlength{\parindent}{0pt}
\setlength{\emergencystretch}{3em}
\setlist{itemsep=0.2em,topsep=0.4em}

% ---------------------------------------------------------------------------
% Code listings
% ---------------------------------------------------------------------------
\lstdefinelanguage{Rust}{
  morekeywords={as,break,const,continue,crate,else,enum,extern,false,fn,for,
    if,impl,in,let,loop,match,mod,move,mut,pub,ref,return,self,Self,static,
    struct,super,trait,true,type,unsafe,use,where,while,dyn,Some,None,Ok,Err,
    Option,Result,Vec,String,usize,isize,u8,u64,bool,char,str},
  sensitive=true,
  morecomment=[l]{//},
  morecomment=[s]{/*}{*/},
  morestring=[b]",
}
\lstdefinelanguage{Toml}{
  morecomment=[l]{\#},
  morestring=[b]",
  sensitive=true,
}
\lstdefinestyle{base}{
  basicstyle=\ttfamily\small\color{ink},
  keywordstyle=\color{keywordc}\bfseries,
  commentstyle=\color{commentc}\itshape,
  stringstyle=\color{stringc},
  columns=fullflexible,
  keepspaces=true,
  showstringspaces=false,
  breaklines=true,
  breakatwhitespace=false,
  postbreak=\mbox{\textcolor{commentc}{$\hookrightarrow$}\space},
  tabsize=4,
  backgroundcolor=\color{codebg},
  frame=single,
  rulecolor=\color{codeframe},
  framesep=4pt,
  xleftmargin=4pt,
  xrightmargin=4pt,
  aboveskip=0.8em,
  belowskip=0.8em,
  upquote=true,
}
\lstdefinestyle{rust}{style=base,language=Rust}
\lstdefinestyle{fortran}{style=base,language=[95]Fortran,morekeywords={contains,module,submodule,associate,block,critical,select,case,class,type,enumeration,procedure,pure,elemental,impure,abstract,interface,extends,only,use,implicit,none,intent,in,out,inout,parameter,allocatable,pointer,target,concurrent},keywordstyle=\color{keywordc}}
\lstdefinestyle{shell}{style=base,language=bash,keywordstyle=\color{ink}}
\lstdefinestyle{plain}{style=base,language={}}
\lstdefinestyle{toml}{style=base,language=Toml}
\lstdefinestyle{before}{style=fortran,backgroundcolor=\color{beforec}}
\lstdefinestyle{after}{style=fortran,backgroundcolor=\color{afterc}}
\lstset{style=base}

\newcommand{\code}[1]{\texttt{\small #1}}
\newcommand{\file}[1]{\texttt{\small\color{accent}#1}}
\newcommand{\fn}[1]{\texttt{\small #1}}
\newcommand{\flag}[1]{\texttt{\small -{}-#1}}

% Side-by-side before/after. Use as
%   \begin{beforeafter} ... input listing ... \splitafter ... output listing ... \end{beforeafter}
% (implemented with two minipages, since listings cannot sit in macro arguments)
\newenvironment{beforecol}{\noindent\begin{minipage}[t]{0.485\linewidth}\textsf{\footnotesize\color{commentc}Input}}{\end{minipage}\hfill}
\newenvironment{aftercol}{\begin{minipage}[t]{0.485\linewidth}\textsf{\footnotesize\color{commentc}Output}}{\end{minipage}\par\medskip}

% ---------------------------------------------------------------------------
% Boxes
% ---------------------------------------------------------------------------
\newcounter{exercise}[chapter]
\renewcommand{\theexercise}{\thechapter.\arabic{exercise}}
\newtcolorbox{exercise}[1][]{enhanced,breakable,colback=exbg,colframe=exframe,
  boxrule=0.6pt,arc=2pt,left=6pt,right=6pt,top=4pt,bottom=4pt,
  fonttitle=\sffamily\bfseries\small,coltitle=white,
  title={Exercise \refstepcounter{exercise}\theexercise\ifx&#1&\else: #1\fi}}
\newtcolorbox{keyidea}[1][Key idea]{enhanced,breakable,colback=accentlight,colframe=accent,
  boxrule=0.6pt,arc=2pt,left=6pt,right=6pt,top=4pt,bottom=4pt,
  fonttitle=\sffamily\bfseries\small,title={#1}}
\newtcolorbox{pitfall}[1][Watch out]{enhanced,breakable,colback=warnbg,colframe=warnframe,
  boxrule=0.6pt,arc=2pt,left=6pt,right=6pt,top=4pt,bottom=4pt,
  fonttitle=\sffamily\bfseries\small,title={#1}}
\newtcolorbox{sourcebox}[1]{enhanced,breakable,colback=white,colframe=codeframe,
  boxrule=0.5pt,arc=1pt,left=6pt,right=6pt,top=3pt,bottom=3pt,
  fonttitle=\ttfamily\small,coltitle=ink,colbacktitle=codebg,title={#1}}

% TikZ styles
\tikzset{
  stage/.style={draw=accent,fill=accentlight,rounded corners=2pt,minimum height=8mm,
    minimum width=30mm,align=center,font=\sffamily\small,inner sep=4pt},
  data/.style={draw=codeframe,fill=codebg,minimum height=7mm,align=center,
    font=\ttfamily\footnotesize,inner sep=3pt},
  note/.style={font=\sffamily\scriptsize,color=commentc,align=left},
  flow/.style={-{Stealth[length=2.2mm]},thick,color=ink!70},
  decision/.style={diamond,draw=exframe,fill=exbg,aspect=2.2,align=center,
    font=\sffamily\scriptsize,inner sep=1pt},
}

\hypersetup{pdftitle={Learn fprettier: A Guide for Developers},pdfauthor={Fred Jones}}

\begin{document}

% ---------------------------------------------------------------------------
% Title page
% ---------------------------------------------------------------------------
\frontmatter
\begin{titlepage}
\centering
\vspace*{3cm}
{\sffamily\color{accent}\large A Guide for Developers\par}
\vspace{0.6cm}
{\sffamily\bfseries\color{ink}\fontsize{42}{48}\selectfont Learn fprettier\par}
\vspace{0.8cm}
{\sffamily\large\color{ink}How a fast Fortran formatter works,\\ from the command line down to the last space\par}
\vspace{2cm}
\begin{tikzpicture}[x=1cm,y=1cm]
  \foreach \i/\w in {0/9,1/7,2/8,3/5,4/6,5/8,6/4}{
    \fill[accent!\the\numexpr 25+\i*10\relax] (\i*0.25,-\i*0.45) rectangle ++(\w*0.6,0.28);
  }
  \foreach \i in {0,...,6}{\fill[codeframe] (-0.6,-\i*0.45) rectangle ++(0.3,0.28);}
\end{tikzpicture}
\vfill
{\sffamily\large Fred Jones\par}
\vspace{0.3cm}
{\sffamily\small\color{commentc} Covers fprettier 0.5.0 and later \textperiodcentered{} November 2026\par}
\end{titlepage}

\thispagestyle{empty}
\vspace*{\fill}
{\small
\textbf{Learn fprettier: A Guide for Developers}\par\medskip
Copyright \copyright{} 2026 Fred Jones.\par\medskip
fprettier is distributed under the terms of the MIT License or the Apache
License 2.0, at your option. This book may be shared and adapted under the same
terms.\par\medskip
The source of this book is the single file \file{book/learn-fprettier.tex} in
the fprettier repository, \url{https://github.com/fredjones42/fprettier}. Build it with
\code{tectonic learn-fprettier.tex}.\par\medskip
Every example showing fprettier output was produced by running fprettier
(version 0.5.0 plus the fixes listed as unreleased in its changelog, as of
November 2026) on the input shown, with default settings unless the text says otherwise.
Source excerpts are taken from the same version; some are abridged, which is
marked where it matters. Function and file names are given so that you can find
the code even after its line numbers have moved.\par}

\tableofcontents

% ---------------------------------------------------------------------------
\chapter{Preface}
% ---------------------------------------------------------------------------

fprettier is an auto-formatter for free-form Fortran (Fortran 90 and later). Give it a
messy source file and it gives back the same program with consistent
indentation, consistent spacing around operators, continuation lines aligned
under the brackets they continue, over-long lines broken, and, if you ask,
keywords in the case you prefer and \code{use} statements in alphabetical order.

It was inspired by \emph{fprettify}, the Python formatter that many Fortran
projects already use, and it keeps fprettify's formatting
conventions closely enough that switching is usually uneventful. It differs in
two ways: it is written in Rust and formats large trees in parallel, and it is
actively maintained, with a test suite of almost four hundred tests that pins
down its behaviour.

\section*{Who this book is for}

The book has two parts, for two kinds of reader.

\textbf{Part I, \emph{Using fprettier},} is for Fortran developers who run the
tool. It explains what fprettier changes and why, how to configure it, how
its indentation rules handle real-world code (including code that was never
consistently indented), and how to keep it away from code you have formatted by
hand. You do not need to know any Rust to read Part I.

\textbf{Part II, \emph{Inside fprettier},} is for developers who want to understand
the implementation, fix a bug in it, or add a feature. It follows a source file
through the program: from the command line, through the character filter that
tells code from strings and comments, the stream that reassembles continuation
lines, the two formatting passes, the scope-tracking indenter, the continuation
aligner, the whitespace engine and the line splitter, to the bytes written
back to disk. It ends with the test suite and a set of recipes for common
changes. Part II assumes you can read Rust; it does not assume you know
Fortran well, and it explains the Fortran rules the code depends on as they come up.

If you are in the second group, read at least Chapters~\ref{ch:tour}
and~\ref{ch:indent-user} of Part I first: Part II explains \emph{how} the code
produces the behaviour that Part I describes.

\section*{How to read the examples}

Most examples are side-by-side pairs, with the input on a pale red background
and fprettier's output on a pale green one:

\begin{beforecol}
\begin{lstlisting}[style=before]
if(x>0)then
y=x**2
endif
\end{lstlisting}
\end{beforecol}
\begin{aftercol}
\begin{lstlisting}[style=after]
if (x > 0) then
    y = x**2
end if
\end{lstlisting}
\end{aftercol}

Every one of them was produced by running fprettier, so the output is exactly
what the tool prints, quirks included. Where the output has a quirk, the text
points it out and explains where it comes from, since tracing a quirk to its
cause is a good way to learn how the formatter works.

Rust excerpts carry the file they come from. Excerpts are sometimes shortened by
dropping comments or unrelated branches; the code that is shown is real.

Each chapter ends with exercises. Part I's exercises ask you to predict
fprettier's output and then check your prediction; Part II's ask you to read,
instrument or change the code. You need a checkout of the repository and a
Rust toolchain (1.87 or later) for most of them:

\begin{lstlisting}[style=shell]
git clone https://github.com/fredjones42/fprettier
cd fprettier
cargo build --release
cargo test
\end{lstlisting}

\mainmatter
% ===========================================================================
\part{Using fprettier}
% ===========================================================================

% ---------------------------------------------------------------------------
\chapter{Getting Started}\label{ch:start}
% ---------------------------------------------------------------------------

\section{Installing}

fprettier is a Rust program distributed on crates.io. With a Rust toolchain
installed (see \url{https://rustup.rs}), install it with Cargo:

\begin{lstlisting}[style=shell]
cargo install fprettier
\end{lstlisting}

The crate needs Rust 1.87 or later. It builds one binary, \code{fprettier}, and
also exposes its formatter as a library (the \code{fprettier} crate), which
Part II uses for experiments.

\section{A first run}

Save this file as \file{geometry.f90}:

\begin{lstlisting}[style=before]
module geometry
implicit none
real,parameter::pi=3.14159
contains
subroutine area(r,a)
real,intent(in)::r
real,intent(out)::a
if(r<0.0)then
a=-1.0
else
a=pi*r**2
endif
end subroutine area
end module geometry
\end{lstlisting}

and run

\begin{lstlisting}[style=shell]
fprettier geometry.f90
\end{lstlisting}

fprettier rewrites the file in place and prints \code{Formatted 1 files
successfully.} on standard error. The file now reads:

\begin{lstlisting}[style=after]
module geometry
    implicit none
    real, parameter::pi = 3.14159
contains
    subroutine area(r, a)
        real, intent(in)::r
        real, intent(out)::a
        if (r < 0.0) then
            a = -1.0
        else
            a = pi*r**2
        end if
    end subroutine area
end module geometry
\end{lstlisting}

Several rules are visible in this one example:

\begin{itemize}
  \item Every scope (\code{module}, \code{subroutine}, \code{if}) indents its
    body by four spaces, the default indent.
  \item \code{contains} and \code{else} sit at the level of the construct they
    belong to, not of its body.
  \item Commas are followed by one space; \code{=} and \code{<} get a space on
    each side; \code{if} is separated from its parenthesis.
  \item \code{endif} is spelled \code{end if}.
  \item \code{pi*r**2} is left tight. At the default whitespace level, 2,
    multiplication and division are not spaced; level 3 would write
    \code{pi * r**2}. Exponentiation (\code{**}) is never spaced.
  \item \code{real, parameter::pi} keeps its \code{::} tight, because
    declarations are left as written unless you pass \flag{format-decl}.
  \item \code{-1.0} stays \code{-1.0}: a unary minus is not spaced.
\end{itemize}

The chapters that follow explain each of these rules and how to change it.

\section{Ways to run fprettier}

\subsection{In place, on files and directories}

Arguments are files or directories. A directory is searched for Fortran files
one level deep; add \flag{recursive} (\code{-r}) to search its whole subtree:

\begin{lstlisting}[style=shell]
fprettier src/main.f90 src/util.f90
fprettier -r src/
\end{lstlisting}

A file inside a directory is formatted when its extension, compared without
regard to case, is one of \code{f90}, \code{f95}, \code{f03}, \code{f08},
\code{f18}, \code{f}, \code{for}, \code{ftn} or \code{fpp}. Add more with
\code{-f}/\flag{fortran}, which may be repeated (\code{-f inc -f F2018}). A
file you name explicitly is always formatted, whatever its extension.

Two more filters apply to what the directory walk finds:
\code{-e}/\flag{exclude} takes a glob pattern (repeatable) that is matched
against the whole path, the file name and each path component, so
\code{-e build} skips everything under any directory called \file{build};
and \code{-m}/\flag{exclude-max-lines}~\textit{N} skips files of more than
\textit{N} lines. Files over 100\,MB are always skipped.

When fprettier writes a file, it writes only if the formatted text differs
from the original, and it writes atomically: the new text goes into a temporary
file beside the original, which is then renamed over it. An interrupted run
cannot leave a half-written source file, the file keeps its permissions, and a
symbolic link keeps pointing at the file it pointed at before.

\subsection{To standard output}

\flag{stdout} (\code{-s}) prints the formatted text instead of writing it back.
This is how every example in this book was produced. With no file arguments,
or with the single argument \code{-}, fprettier reads standard input and
writes standard output, which is what editor integrations use:

\begin{lstlisting}[style=shell]
fprettier < messy.f90 > tidy.f90
\end{lstlisting}

If standard input is a terminal and there are no arguments, fprettier prints
its help instead of waiting for input.

\subsection{In CI: \texttt{-{}-check} and \texttt{-{}-diff}}

\flag{check} formats in memory, changes nothing, prints
\code{Would reformat: \textit{path}} for every file whose formatting would
change, and exits with status 1 if there were any. \flag{diff} (\code{-d}) prints a
unified diff of the changes instead. They combine:

\begin{lstlisting}[style=shell]
$ fprettier --check --diff f.f90
Would reformat: f.f90
--- f.f90
+++ f.f90
@@ -1,4 +1,4 @@
 program p
-integer::i
-i=1
+    integer::i
+    i = 1
 end program p
$ echo $?
1
\end{lstlisting}

A CI job needs only \code{fprettier --check -r src/}. Both modes are safe to
pipe into \code{head} or \code{less}: when the reader closes the pipe early,
fprettier exits quietly instead of reporting a broken pipe.

\subsection{Exit status}

fprettier exits with 0 on success and 1 if any file could not be formatted,
if an input path does not exist (it reports \code{no such file or directory}
rather than silently formatting nothing), if the configuration is invalid, or,
under \flag{check}, if any file would change.

\subsection{Parallelism}

When it is writing files in place, fprettier formats them in parallel across
all cores. \code{-j}/\flag{jobs}~\textit{N} sets the number of threads;
\code{-j 1} formats one file at a time. \flag{stdout}, \flag{diff} and
\flag{check} always run sequentially, so that their output comes out in a
stable order.

\subsection{Quiet and verbose}

\code{-S}/\flag{silent} suppresses the summary line and the skip messages, for
editor integrations. \code{-D}/\flag{debug} prints the configuration files
found, the final configuration, and any in-file directive, which is the first
thing to try when fprettier does not do what you expect
(Section~\ref{sec:debug}).

\begin{exercise}[Your own tree]
Run \code{fprettier --check -r} on a Fortran project you work on, then run it
again with \flag{diff} and page through the result. Which kinds of change are
the most common? Keep the list: by the end of Part~I you should be able to say,
for each kind of change, which setting controls it.
\end{exercise}

\begin{exercise}[Exit codes]
Run \code{fprettier --check nosuchfile.f90; echo \$?}. Then create an empty
directory and run \code{fprettier --check emptydir; echo \$?}. Explain the
difference between the two results.
\end{exercise}

% ---------------------------------------------------------------------------
\chapter{A Tour of What fprettier Changes}\label{ch:tour}
% ---------------------------------------------------------------------------

This chapter walks through each kind of change fprettier makes, with a small
example of each. Later chapters give the full rules; the goal here is to know
what to expect.

\section{Indentation}

Every construct that has a body indents it by one level: program units
(\code{program}, \code{module}, \code{submodule}), procedures
(\code{subroutine}, \code{function}), \code{interface} blocks, derived-type
and \code{enum} definitions, and the block constructs \code{if}, \code{do},
\code{select case}/\code{type}/\code{rank}, \code{where}, \code{forall},
\code{associate}, \code{block}, \code{critical} and \code{change team}.
Statements that divide a construct, such as \code{else}, \code{else if},
\code{case}, \code{contains} and \code{elsewhere}, go at the level of the
construct itself.

\begin{beforecol}
\begin{lstlisting}[style=before]
subroutine classify(n)
integer::n
select case(n)
case(:-1)
print*,"negative"
case(0)
print*,"zero"
case default
print*,"positive"
end select
end subroutine
\end{lstlisting}
\end{beforecol}
\begin{aftercol}
\begin{lstlisting}[style=after]
subroutine classify(n)
    integer::n
    select case (n)
    case (:-1)
        print *, "negative"
    case (0)
        print *, "zero"
    case default
        print *, "positive"
    end select
end subroutine
\end{lstlisting}
\end{aftercol}

Note that \code{case} lines are \emph{not} indented relative to
\code{select case}. That is fprettify's convention, and fprettier keeps it.

Indentation is the most subtle part of fprettier, because it also has to cope
with code that was indented inconsistently, or deliberately left unindented, in
the first place. Chapter~\ref{ch:indent-user} covers it in detail.

\section{Whitespace}

fprettier first removes the spacing a statement has, apart from single
blanks between words, and then adds back spacing according to a set of eleven
rules: after commas, around assignments, around relational and logical
operators, and so on. A \emph{whitespace level} from 0 to 4 turns groups of rules
on, and each rule can also be switched on or off on its own:

\begin{beforecol}
\begin{lstlisting}[style=before]
x  =  a( i , j )+b  *c
if(x.gt.0.and.y/=1)call  foo (x)
call s(n=1, m=2)
y=-x+val_1d-2
loop:do i=1,3
\end{lstlisting}
\end{beforecol}
\begin{aftercol}
\begin{lstlisting}[style=after]
x = a(i, j) + b*c
if (x .gt. 0 .and. y /= 1) call foo(x)
call s(n=1, m=2)
y = -x + val_1d - 2
loop: do i = 1, 3
\end{lstlisting}
\end{aftercol}

Note a few details. The \code{=} of a keyword argument (\code{n=1}) is not
spaced, because it is inside parentheses. The \code{-2} after \code{val\_1d} is
recognised as a binary minus even though \code{1d-2} looks like the start of a
double-precision literal, because \code{val\_1d} begins with a letter. A
construct name gets a space after its colon. Chapter~\ref{ch:whitespace-user}
lists every rule.

\section{Continuation alignment}

A statement continued with \code{\&} has its continuation lines aligned to a
column that depends on where the line was broken: just inside the innermost
open bracket, after the \code{=} of an assignment, or after the \code{::} of a
declaration. If none of these applies, the continuation is indented by one
level.

\begin{beforecol}
\begin{lstlisting}[style=before]
call compute(alpha, beta, &
gamma, delta)
total = first_value + &
second_value
integer, dimension(3) :: a, &
b
call foo( &
a, b)
\end{lstlisting}
\end{beforecol}
\begin{aftercol}
\begin{lstlisting}[style=after]
call compute(alpha, beta, &
             gamma, delta)
total = first_value + &
        second_value
integer, dimension(3) :: a, &
                         b
call foo( &
    a, b)
\end{lstlisting}
\end{aftercol}

The last example shows an exception: a bracket that is the last thing on its
line (\code{foo( \&}) does not set the alignment column, since aligning after
it would waste the line; the continuation is indented by one level instead.

\section{Long lines}

Lines longer than the line-length limit (132 by default, \flag{line-length}
or \code{-l} to change it) are broken at a space, or failing that after a comma,
and the pieces are aligned by the rules above. An inline comment that no longer
fits beside its code moves to a line of its own below it:

\begin{lstlisting}[style=before]
subroutine s
    call long_procedure_name(argument_number_one, argument_number_two, argument_number_three, arg4)
    x = 1 ! a fairly long trailing comment that pushes the line past sixty
end subroutine
\end{lstlisting}

With \code{-l 60}:

\begin{lstlisting}[style=after]
subroutine s
    call long_procedure_name(argument_number_one, &
                             argument_number_two, &
                             argument_number_three, arg4)
    x = 1
    ! a fairly long trailing comment that pushes the line past sixty
end subroutine
\end{lstlisting}

fprettier does not break inside a string, and it does not make a break that
would not help: if the only way to break a line leaves it as long as before, the
line is left alone. Section~\ref{sec:linesplit} explains the details.

\section{Case conversion}

Off by default. \flag{case} takes four values, for keywords, intrinsic
procedures, logical and relational operators, and named constants from
\code{iso\_fortran\_env} and \code{iso\_c\_binding}; each is 0 (leave alone),
1 (lowercase) or 2 (uppercase). \flag{case-types} sets a fifth category, kind
values and exponent letters:

\begin{lstlisting}[style=before]
x=1.0e-3_dp
real(kind=dp)::a=2.5e-8_real64
if(a.EQ.b.AND.c)then
n=size(a)
size=3
end if
\end{lstlisting}

With \code{-{}-case 1 1 1 1 -{}-case-types 2}:

\begin{lstlisting}[style=after]
x = 1.0E-3_DP
real(kind=DP)::a = 2.5E-8_REAL64
if (a .eq. b .and. c) then
    n = size(a)
    size = 3
end if
\end{lstlisting}

Intrinsic procedure names are converted only where they are \emph{called}:
\code{size(a)} is a call, but the variable called \code{size} on the next line
is not touched. Your own identifiers are never touched.

\section{Operator replacement}

Off by default. \flag{enable-replacements} rewrites the relational operators to one
spelling: Fortran-style (\code{.lt.}, \code{.eq.}, \ldots) by default, or
C-style (\code{<}, \code{==}, \ldots) with \flag{c-relations}. With both flags,
\code{y=x .lt. 3 .and. .not. flag} becomes \code{y = x < 3 .and. .not. flag}.
Logical operators have no C-style spelling in Fortran and are left alone.

\section{Sorting \texttt{use} statements}

Off by default. \flag{sort-use} sorts each run of consecutive \code{use}
statements by module name; \flag{sort-use-only} sorts the names in each
\code{only:} list. A rename such as \code{a => q} sorts by its local name,
\code{a}:

\begin{beforecol}
\begin{lstlisting}[style=before]
module m
use zeta
use alpha, only: z, b, a => q
\end{lstlisting}
\end{beforecol}
\begin{aftercol}
\begin{lstlisting}[style=after]
module m
    use alpha, only: a => q, b, z
    use zeta
\end{lstlisting}
\end{aftercol}

\section{Small normalisations}

A few more changes happen without any option:

\begin{itemize}
  \item A run of blank lines collapses to a single blank line.
  \item Trailing whitespace is removed.
  \item Tabs are expanded to eight spaces, except inside character literals,
    where a tab is part of the value.
  \item \code{endif}, \code{enddo}, \code{endselect} and the other run-together
    \code{end}s become \code{end if}, \code{end do}, \ldots{} (but a variable
    called \code{endif}, as in \code{endif = 3}, is left alone).
  \item A statement label is followed by exactly enough space to put the
    statement at its indentation column.
  \item Windows (CRLF) line endings are preserved.
\end{itemize}

fprettier also warns, on standard error, about input that is already over the
limits the Fortran 2023 standard sets for free source form: 10\,000 characters
per line and 1\,000\,000 characters per statement. It formats such a file
anyway.

\section{What fprettier will not do}

fprettier is a formatter, not a compiler. It does not parse Fortran into a
syntax tree; it recognises statements with regular expressions and tracks
brackets, strings and comments with a small state machine. In practice this means:

\begin{itemize}
  \item It never changes what a program means. Strings, comments and
    preprocessor lines are reproduced exactly; only blanks outside them are
    changed, along with the case changes and operator spellings you ask for.
  \item It does not reject invalid code. Unbalanced brackets or an \code{end
    do} without a \code{do} are handled as well as possible, and the rest of the
    file is still formatted.
  \item It handles only free-form source. Fixed-form files (columns 1--6
    reserved, \code{C} in column 1 for comments) are not supported, even though
    \code{.f} and \code{.for} are in the default extension list.
\end{itemize}

\begin{keyidea}[Idempotence]
Formatting a formatted file changes nothing. fprettier is designed and tested
so that its output is a fixed point: \code{fprettier --check} passes on any file
fprettier has just written. Several of the more intricate pieces of code in
Part II exist only to make this true.
\end{keyidea}

\begin{exercise}[Predict, then check]
Before running it, write down what fprettier will do to each line:
\begin{lstlisting}[style=fortran]
print*,'a'//'b'
do while(i<10)
x(1:n)=y(1:n)*2
p=>t
\end{lstlisting}
Then check with \code{fprettier -s}. Which of your predictions were wrong?
\end{exercise}

\begin{exercise}[Case conversion boundaries]
Format \code{integer :: size; size = size(arr)} with \code{-{}-case 2 2 2 2}.
Which occurrences of \code{size} change case and which do not? Why can the
formatter tell them apart?
\end{exercise}

% ---------------------------------------------------------------------------
\chapter{Configuration}\label{ch:config}
% ---------------------------------------------------------------------------

\section{Three layers}

fprettier takes its settings from three places. Each overrides the one before it:

\begin{enumerate}
  \item \file{fprettier.toml} configuration files;
  \item command-line flags;
  \item an \code{! fprettier:} directive comment in the file being formatted.
\end{enumerate}

\begin{figure}[ht]
\centering
\begin{tikzpicture}[node distance=3mm]
  \node[data,minimum width=40mm] (def) {Config::default()};
  \node[data,minimum width=40mm,below=of def] (home) {\textasciitilde/fprettier.toml};
  \node[data,minimum width=40mm,below=of home] (root) {/fprettier.toml};
  \node[data,minimum width=40mm,below=of root] (mid) {/work/fprettier.toml};
  \node[data,minimum width=40mm,below=of mid] (leaf) {/work/src/fprettier.toml};
  \node[data,minimum width=40mm,below=of leaf] (cli) {command-line flags};
  \node[data,minimum width=40mm,below=of cli] (dir) {! fprettier: directive};
  \node[stage,right=22mm of mid,minimum width=24mm] (cfg) {effective\\Config};
  \foreach \n in {def,home,root,mid,leaf,cli,dir}{\draw[flow] (\n.east) -- ++(8mm,0) |- (cfg.west);}
  \draw[-{Stealth[length=2.2mm]},thick,commentc] ([xshift=-8mm]def.west) -- node[note,sloped,above]{later overrides earlier} ([xshift=-8mm]dir.west);
\end{tikzpicture}
\caption{Configuration layers for a file in \file{/work/src}. Later layers
override earlier ones, key by key.}
\label{fig:config-layers}
\end{figure}

\section{Configuration files}

\subsection{Where fprettier looks}

For each file it formats, fprettier looks for \file{fprettier.toml} in your
home directory, then in every directory from the filesystem root down to the
directory holding the file. Files found later win, so a \file{fprettier.toml}
beside the source overrides one at the top of the repository, which overrides
your personal one. Only the keys a file actually sets are applied, and nested
tables (\code{[whitespace\_dict]} and \code{[case\_dict]}) are merged key by
key, so a subdirectory can change one whitespace rule without repeating the
others.

\code{-c}/\flag{config}~\textit{PATH} uses exactly that one file and skips the
search. An empty file there is a convenient way to say ``defaults only'':

\begin{lstlisting}[style=shell]
touch /tmp/empty.toml
fprettier --config /tmp/empty.toml -r src/
\end{lstlisting}

\subsection{Keys}

Every key is optional. This file lists them all with their defaults:

\begin{lstlisting}[style=toml]
indent = 4                      # spaces per indent level
line_length = 132               # maximum line length (0: no limit)
whitespace = 2                  # whitespace level, 0-4
impose_indent = true            # apply indentation
impose_whitespace = true        # apply whitespace formatting
strict_indent = false           # re-indent IF/DO even when already aligned
indent_fypp = true              # indent fypp preprocessor blocks
indent_mod = true               # indent module/program/submodule bodies
format_decl = false             # normalise spacing around ::
comment_spacing = 1             # spaces before an inline comment
normalize_comment_spacing = false   # apply comment_spacing everywhere
enable_replacements = false     # convert between .lt. and < style operators
c_relations = false             # with enable_replacements, prefer < over .lt.
sort_use = false                # sort use statements within a group
sort_use_only = false           # sort the names in a use ... only: list

[whitespace_dict]               # per-rule overrides of the level
# comma = true, multdiv = true, ...   (see Chapter 4)

[case_dict]                     # 0 = leave, 1 = lower, 2 = upper
# keywords = 1, procedures = 1, operators = 1, constants = 1, types = 2
\end{lstlisting}

An unknown key is an error, reported against the file that contains it, as is a
file that is not valid TOML. fprettier refuses to carry on with defaults in that
case, because silently reformatting a whole tree with settings nobody chose is
worse than stopping.

\section{Command-line flags}

Every setting has a flag; Appendix~\ref{app:cli} lists them. Three things about
how flags combine with the configuration file are worth knowing.

\textbf{Value flags replace.} \code{-i}, \code{-l}, \code{-w} and
\flag{comment-spacing} replace the configured value.

\textbf{Boolean flags only move one way.} A \flag{no-*} flag can only switch
off a setting the configuration turned on (\flag{no-indent},
\flag{no-whitespace}, \flag{no-indent-fypp}, \flag{no-indent-mod}), and the
other boolean flags (\flag{strict-indent}, \flag{format-decl},
\flag{sort-use}, \ldots) can only switch one on. There is no flag to switch
\code{sort\_use} back off if a configuration file turned it on; use
\flag{config} with a different file for that.

\textbf{Whitespace rule flags take a value.} Each of the eleven rules has a
flag \code{-{}-whitespace-\textit{rule}}, which takes an optional
\code{=true} or \code{=false}; the bare flag means true:

\begin{lstlisting}[style=shell]
fprettier -w 2 --whitespace-multdiv --whitespace-plusminus=false file.f90
\end{lstlisting}

\section{In-file directives}

A comment of the form \code{! fprettier: \textit{options}} anywhere in a file
overrides everything else for that file:

\begin{lstlisting}[style=fortran]
! fprettier: --indent 2 --line-length 100 --case 1 1 1 1 --case-types 2
\end{lstlisting}

Directives accept a subset of the flags: \code{-i}/\flag{indent},
\code{-l}/\flag{line-length}, \code{-w}/\flag{whitespace}, \flag{case},
\flag{case-types}, \flag{no-indent}, \flag{enable-indent},\linebreak[2]
\flag{no-whitespace} and \flag{enable-whitespace}. Anything else is ignored.
Only the \emph{first} directive in a file is read; put every option on that
one line. Unlike a command-line \flag{no-*} flag, a directive can switch
indentation or whitespace formatting back on with \flag{enable-indent} or
\flag{enable-whitespace}.

The directive parser is lenient: an option it does not recognise, or a value it
cannot parse, is skipped. The resulting settings are then validated like the
command line's, so \code{! fprettier: -{}-indent 0} is an error.

\section{Validation}

Only two values are checked: \code{indent} must be at least 1, and
\code{line\_length} must be at most 10\,000, the longest line free source form
permits. Other values degrade harmlessly: a whitespace level above 4 behaves
like level 2, and a line length below 40 effectively turns line splitting off.
A \code{line\_length} of 0 means ``no limit''.

\section{Seeing what fprettier decided}\label{sec:debug}

\flag{debug} prints where the configuration came from and what it ended up as:

\begin{lstlisting}[style=plain]
$ fprettier -s -D --config /dev/null f.f90
[DEBUG] Using explicit config file: /dev/null
[DEBUG] Configuration: Config {
    indent: 4,
    line_length: 132,
    whitespace: 2,
    whitespace_dict: {},
    impose_indent: true,
    ...
}
[DEBUG] whitespace_flags array: [true, true, true, true, true, false,
                                 true, false, true, true, false]
\end{lstlisting}

Without \flag{config} it lists every \file{fprettier.toml} it discovered, in
order. The \code{whitespace\_flags} array is the resolved state of the eleven
whitespace rules; Table~\ref{tab:ws-levels} on page~\pageref{tab:ws-levels}
explains its positions.

\begin{exercise}[Layering]
Create \file{a/fprettier.toml} with \code{indent = 2} and
\code{[whitespace\_dict] multdiv = true}, and \file{a/b/fprettier.toml} with
\code{indent = 3}. Format a file in \file{a/b} with \flag{debug}. Which
indent and which whitespace flags are used? Now add \code{-i 5} on the command
line, and then \code{! fprettier: -i 6} inside the file.
\end{exercise}

\begin{exercise}[A directive that does nothing]
Put \code{! fprettier: -{}-sort-use} in a file and format it. Was anything
sorted? Explain using the list of options directives accept.
\end{exercise}
% ---------------------------------------------------------------------------
\chapter{Whitespace Rules}\label{ch:whitespace-user}
% ---------------------------------------------------------------------------

\section{Levels and rules}

Whitespace formatting is controlled by eleven independent rules. The
\code{whitespace} setting (\code{-w}) picks a preset that switches groups of
them on; \code{[whitespace\_dict]} in a configuration file, or the
\code{-{}-whitespace-\textit{rule}} flags, then override single rules.
Table~\ref{tab:ws-levels} shows which rules each level enables. The
\emph{index} column is the rule's position in the \code{whitespace\_flags}
array that \flag{debug} prints.

\begin{table}[ht]
\centering\small
\begin{tabularx}{\linewidth}{@{}c l X ccccc@{}}
\toprule
Index & Key & What it spaces & 0 & 1 & 2 & 3 & 4\\
\midrule
0 & \code{comma}       & after \code{,} and \code{;}; after \code{only:} in \code{use} & & \checkmark & \checkmark & \checkmark & \checkmark\\
1 & \code{assignments} & around \code{=} and \code{=>} (outside brackets) & & \checkmark & \checkmark & \checkmark & \checkmark\\
2 & \code{relational}  & around \code{<} \code{==} \code{/=} \code{.eq.} \ldots, and F2023 \code{?}/\code{:} & & \checkmark & \checkmark & \checkmark & \checkmark\\
3 & \code{logical}     & around \code{.and.} \code{.or.} \code{.not.} \ldots{} and user-defined operators & & \checkmark & \checkmark & \checkmark & \checkmark\\
4 & \code{plusminus}   & around binary \code{+} and \code{-} & & & \checkmark & \checkmark & \checkmark\\
5 & \code{multdiv}     & around binary \code{*} and \code{/} & & & & \checkmark & \checkmark\\
6 & \code{print}       & \code{print*} $\to$ \code{print *} & & \checkmark & \checkmark & \checkmark & \checkmark\\
7 & \code{type}        & around the component selector \code{\%} & & & & & \checkmark\\
8 & \code{intrinsics}   & before a statement keyword's \code{(}; \code{endif} $\to$ \code{end if}; \code{name:do} $\to$ \code{name: do} & & \checkmark & \checkmark & \checkmark & \checkmark\\
9 & \code{decl}        & around \code{::} (only with \flag{format-decl}) & & \checkmark & \checkmark & \checkmark & \checkmark\\
10 & \code{concat}     & around \code{//} & & & & & \checkmark\\
\bottomrule
\end{tabularx}
\caption{Whitespace rules and the levels that enable them. Level 2 is the
default; a level above 4 behaves like level 2.}
\label{tab:ws-levels}
\end{table}

The same line at three levels:

\begin{lstlisting}[style=plain]
input    x  =  a( i , j )+b  *c
-w 0     x=a(i,j)+b*c
-w 2     x = a(i, j) + b*c
-w 3     x = a(i, j) + b * c
\end{lstlisting}

The level-0 output is not simply the input: even with every rule off,
fprettier removes runs of blanks and blanks before punctuation.
Section~\ref{sec:ws-stages} follows this line through the whitespace engine.

\section{What each rule does, precisely}

\paragraph{Commas and semicolons.} No space before, exactly one after. The
\code{only:} of a \code{use} statement is also followed by one space.

\paragraph{Assignments.} \code{=} gets one space on each side, but only at
bracket depth zero, so keyword arguments (\code{call s(n=1)}), and
\code inside \code{real(\ldots)}, are left tight. A pointer
assignment \code{=>} is spaced at any depth, which covers the
associations of \code{associate (y => x)}. The \code{=} inside \code{==},
\code{/=}, \code{<=} and \code{>=} is not an assignment and is left to the
relational rule.

\paragraph{Relational operators.} All twelve spellings are spaced on both
sides. The \code{?} and \code{:} of a Fortran 2023 conditional expression,
\code{(a > 0 ? 1 : 2)}, are spaced the same way. fprettier tells such a colon apart from an
array section's (\code{a(1:n)}) or a \code{case (1:2)}: a colon is spaced only
inside a bracket that already contains a \code{?}.

\paragraph{Logical and user-defined operators.} \code{.and.}, \code{.or.},
\code{.not.}, \code{.eqv.} and \code{.neqv.} are spaced, and so is any
user-defined operator of the form \code{.\textit{name}.}, so
\code{x = a.cross.b} becomes \code{x = a .cross. b}. The logical literals
\code{.true.} and \code{.false.} have the same shape but are left alone,
including any kind suffix (\code{.true.\_lk}).

\paragraph{Plus and minus.} Only binary \code{+} and \code{-} are spaced. A
sign is binary if what comes before it is a name, a number, \code{)} or
\code{]}; otherwise it is unary and stays attached (\code{-1.0},
\code{(-x)}, \code{a*(-b)}). The sign of an exponent (\code{1.0d-3},
\code{2e+5}) is never spaced; fprettier checks that the text before the
\code{d} or \code{e} really is a number, so in \code{val\_1d-2} the minus is
binary. At the start of a continuation line, fprettier looks at the last
character of the previous line to decide.

\paragraph{Multiply and divide.} Binary \code{*} and \code{/} are spaced at
level 3 and above. \code{**} is never spaced, and neither is the \code{/} of
an array constructor's \code{(/} and \code{/)}.

\paragraph{Print and read.} \code{print*,x} becomes \code{print *, x}.

\paragraph{Component selector.} At level 4, \code{a\%b} becomes \code{a \% b}.
At lower levels any blanks around \code{\%} are removed.

\paragraph{Intrinsics.} Despite its name, this rule is about statement
keywords. It puts one space between a statement keyword and its parenthesis:
\code{if}, \code{else if}, \code{do while}, \code{select case}, \code{case},
\code{select type}, \code{type is}, \code{class is}, \code{allocate},
\code{deallocate}, \code{open}, \code{close}, \code{read}, \code{write},
\code{inquire}, \code{forall}, \code{where}, \code{associate} and a few others.
It does not apply to a procedure called \code{open} through a component
(\code{obj\%open(3)}), or to \code{read}/\code{write} naming a defined-I/O generic
(\code{generic :: read(formatted) => \ldots}). The same rule also separates run-together
\code{end} keywords (\code{endif} $\to$ \code{end if}, including after a
semicolon) and puts a space after a construct name's colon (\code{outer:do}
$\to$ \code{outer: do}). Function calls are never separated from their
parentheses: \code{sqrt (x)} becomes \code{sqrt(x)}.

\paragraph{Declarations.} By default fprettier preserves the spacing around
\code{::}, because many codebases align their declarations in columns by hand:

\begin{lstlisting}[style=fortran]
integer, intent(in)                                :: r, i, j, k
integer                                            :: l
\end{lstlisting}

With \flag{format-decl}, \code{::} gets exactly one space on each side and
such columns are collapsed. The \code{decl} rule has an effect only when
\flag{format-decl} is given.

\paragraph{Concatenation.} At level 4, \code{//} is spaced; below that, blanks
around it are removed (\code{'a' // 'b'} $\to$ \code{'a'//'b'}).

\section{Brackets}

Blanks just inside brackets are removed: \code{a( i , j )} $\to$
\code{a(i, j)}, \code{[ 1, 2 ]} $\to$ \code{[1, 2]}. A space is put before an
opening bracket after a closing one or after a comma
(\code{write(*,*)(merge(\ldots))} $\to$ \code{write (*, *) (merge(\ldots))}),
after a top-level assignment (\code{arr =[1]} $\to$ \code{arr = [1]}), and
before an array constructor that follows an operator.

\section{What is never touched}

The inside of a character literal, the text of a comment, a fypp inline
expression (\code{\$\{\ldots\}\$}, \code{\#\{\ldots\}\#}, \code{@\{\ldots\}@}),
a C preprocessor line (\code{\#ifdef}, \code{\#include}, \ldots) and a fypp
directive line (\code{\#:if}, \code{\$:}, \ldots) are reproduced exactly.

\begin{exercise}[Find the level]
Your project's style puts spaces around \code{*} but not around \code{//},
and never around \code{\%}. Which level and which overrides express that? Write
it both as a \file{fprettier.toml} and as command-line flags.
\end{exercise}

\begin{exercise}[Unary or binary?]
Predict the output of
\code{y=a*-b+c(-1)-2e-3-x\%e-3} at level 3, then check. Which rule
decided each sign?
\end{exercise}

% ---------------------------------------------------------------------------
\chapter{How Indentation Works}\label{ch:indent-user}
% ---------------------------------------------------------------------------

\section{Scopes}

fprettier indents by tracking \emph{scopes}. A statement such as
\code{if (\ldots) then} or \code{subroutine s} \emph{opens} a scope; its body
is indented one level deeper than the statement; and the matching \code{end}
\emph{closes} the scope and returns to the statement's level. Some statements
\emph{continue} a scope: \code{else}, \code{case}, \code{contains} and the
like are written at the level of the opening statement, and the body after
them is indented again.

Appendix~\ref{app:scopes} lists every scope fprettier knows. Three details
deserve mention here.

\begin{itemize}
  \item A bare \code{end} closes whatever scope is innermost, as in Fortran 77
    style \code{end} for a subroutine.
  \item \code{where (mask)} and \code{forall (\ldots)} open a scope only when
    nothing follows the parenthesis. \code{where (x > 0) y = 1} is a single
    statement and does not.
  \item An \code{end} that closes nothing fprettier recognised (because the
    opening statement used a form it does not match) is ignored for indentation
    purposes, so one unrecognised construct costs one mis-indented block, not
    the rest of the file.
\end{itemize}

\section{Program units: \texttt{indent\_mod}}

By default, the bodies of \code{program}, \code{module} and \code{submodule}
are indented like any other scope. Many projects prefer to leave them at column
0; \flag{no-indent-mod} (or \code{indent\_mod = false}) does that:

\begin{beforecol}
\begin{lstlisting}[style=after]
program p
    implicit none
    do i = 1, n
        x = x + i
    end do
end program
\end{lstlisting}
\end{beforecol}
\begin{aftercol}
\begin{lstlisting}[style=after]
program p
implicit none
do i = 1, n
    x = x + i
end do
end program
\end{lstlisting}
\end{aftercol}

(Here both sides are fprettier output: on the left with the default, on the
right with \flag{no-indent-mod}.) Procedures, interfaces and type definitions
are always indented.

\section{Respecting existing style: stacked IF and DO}\label{sec:stacking}

A popular way to write tightly nested loops is to keep them in one column:

\begin{lstlisting}[style=fortran]
    do i = 1, n
    do j = 1, m
        a(i,j) = 0
    end do
    end do
\end{lstlisting}

fprettier, like fprettify, preserves this. Its rule is:

\begin{keyidea}[The stacking rule]
When the first statement in a scope's body is an \code{if \ldots{} then} or a
\code{do}, and it was written in the \emph{same column} as the statement that
opened the scope, then the body of that scope is not indented. The two
statements stay \emph{stacked}.
\end{keyidea}

Three conditions switch stacking off, and the body is indented as usual:

\begin{enumerate}
  \item The column is not a multiple of the indent width. Code indented with
    three spaces, formatted with the default four, is re-indented throughout;
    stacking is only kept where it is consistent with the new indent.
  \item The opening statement is a construct whose body is always indented:
    \code{associate}, \code{select}, \code{block}, \code{critical},
    \code{change team}, \code{where}, \code{forall}, \code{interface}, a type
    or \code{enum} definition. An \code{if} written flush against an
    \code{associate} is a new level, not a stack.
  \item \flag{strict-indent} is given. It switches the rule off everywhere.
\end{enumerate}

The opening statement \emph{can} be a program unit or procedure header. This
has an effect that surprises people the first time they see it. Take a file
that was never indented at all:

\begin{beforecol}
\begin{lstlisting}[style=before]
program p
do i=1,n
if(a(i)>0)then
b(i)=1
end if
end do
y=x
end program
\end{lstlisting}
\end{beforecol}
\begin{aftercol}
\begin{lstlisting}[style=after]
program p
do i = 1, n
if (a(i) > 0) then
    b(i) = 1
end if
end do
y = x
end program
\end{lstlisting}
\end{aftercol}

The \code{do} is the first statement of the program and is in the program's
column, so the program body is stacked; the \code{if} is the first statement of
the loop and in its column, so the loop body is stacked; only \code{b(i) = 1},
which is not an \code{if} or \code{do}, gets indented. Formatting with
\flag{strict-indent} gives the fully indented layout:

\begin{lstlisting}[style=after]
program p
    do i = 1, n
        if (a(i) > 0) then
            b(i) = 1
        end if
    end do
    y = x
end program
\end{lstlisting}

\begin{pitfall}[Formatting unindented code]
If you are formatting a codebase for the first time and some of it has no
indentation at all, run fprettier once with \flag{strict-indent}. The
stacking rule is there to \emph{preserve} a style someone chose; it cannot tell
that choice from no choice. After the first run, the default rule keeps
the result stable.
\end{pitfall}

\section{Statement labels}

A numeric statement label is written in column 1, and the statement after it is
placed at its proper indentation column, padded with as many blanks as needed
(at least one):

\begin{beforecol}
\begin{lstlisting}[style=before]
subroutine s
x=1
  100   format(a, &
      i3)
end subroutine
\end{lstlisting}
\end{beforecol}
\begin{aftercol}
\begin{lstlisting}[style=after]
subroutine s
    x = 1
100 format(a, &
           i3)
end subroutine
\end{lstlisting}
\end{aftercol}

A labelled \code{do 100 i = 1, 10} loop is closed by the statement that carries
label 100, the way the Fortran standard defines it, including when several
loops share one terminating statement.

\section{Preprocessors}

\subsection{C preprocessor}

Lines starting with \code{\#} followed by anything but \code{!}, \code{:},
\code{\{} or \code{\}} are C preprocessor directives (\code{\#ifdef},
\code{\#include}, \code{\#define}, \ldots). They are moved to column 0, which
is where most C preprocessors require them, and otherwise left exactly as they
are. They do not affect scope tracking: an \code{\#ifdef} does not indent the
code inside it.

\subsection{fypp}

fypp is a Python-based preprocessor popular in Fortran projects. Its block
directives (\code{\#:if}/\code{\#:elif}/\code{\#:else}/\code{\#:endif},
\code{\#:for}, \code{\#:def}, \code{\#:block}, \code{\#:call},
\code{\#:mute}) open and close scopes just like Fortran constructs, and the
directive lines themselves are indented with the code:

\begin{beforecol}
\begin{lstlisting}[style=before]
subroutine s(x)
real :: x
#:for k in [4, 8]
x = x + ${k}$
#:endfor
end subroutine
\end{lstlisting}
\end{beforecol}
\begin{aftercol}
\begin{lstlisting}[style=after]
subroutine s(x)
    real :: x
    #:for k in [4, 8]
        x = x + ${k}$
    #:endfor
end subroutine
\end{lstlisting}
\end{aftercol}

\flag{no-indent-fypp} turns this off: directive lines keep the indentation they
had, and do not open scopes. Inline fypp expressions
(\code{\$\{k\}\$}) are copied verbatim wherever they appear, and fypp lines are
never case-converted, because fypp recognises its own keywords only in lower
case.

\subsection{OpenMP}

OpenMP directives (\code{!\$omp parallel do} and friends) are comments as far as
Fortran is concerned, and fprettier leaves them in column 0. OpenMP
\emph{conditional compilation} lines, which start with \code{!\$} followed by a space,
contain code that is compiled only with OpenMP enabled. fprettier formats that
code like any other, keeps the \code{!\$} sentinel in column 0, and pads after it
so that the code lands at its indentation column:

\begin{lstlisting}[style=after]
!$  nthreads = omp_get_num_threads()
!$OMP PARALLEL DO
    do i = 1, n
\end{lstlisting}

\section{Comments}

A comment on a line of its own is indented with the code around it. A comment
after code (an \emph{inline} comment) keeps the gap that separated it from the
code, unless \flag{normalize-comment-spacing} is given, in which case the gap
becomes \flag{comment-spacing} spaces (default 1). If an inline comment would
push its line past the length limit, the comment moves to its own line below the
code.

Comments that begin with \code{!!}, the convention of the FORD documentation
generator, keep the indentation they were written with, because FORD comments
are often aligned with the declarations they document.

\begin{exercise}[Stacking]
Format the stacked double loop at the start of Section~\ref{sec:stacking},
indented by four spaces inside a subroutine. Then indent it by three spaces
instead and format it again. Explain the difference using the three
conditions.
\end{exercise}

\begin{exercise}[A blank line between]
Does a blank line between \code{program p} and a following \code{do} change
whether they stack? Form a hypothesis, then test it. (Part II,
Section~\ref{sec:inspect}, explains the answer.)
\end{exercise}

% ---------------------------------------------------------------------------
\chapter{Escape Hatches}\label{ch:escape}
% ---------------------------------------------------------------------------

No formatter gets everything right. fprettier gives you three ways to keep it
away from code you have laid out by hand.

\section{Deactivating formatting: \texttt{!\&}, \texttt{!\&<}, \texttt{!\&>}}

A comment that starts with \code{!\&} at the end of a statement leaves that
statement (all of its continuation lines) exactly as written. \code{!\&<} and
\code{!\&>} mark a block to leave alone; both of the lines that carry the
markers are part of the block:

\begin{lstlisting}[style=after]
  subroutine s()
matrix = [1, 0, &   !&<
          0, 1]     !&>
y=1;z=2  !&
    x = [1, 2, 3, &
        & 4, 5, 6]
\end{lstlisting}

The two deactivated statements are not even re-indented. This is fprettier
output: the input had the same first three lines.

\section{Manual alignment with a leading \texttt{\&}}

Fortran allows a continuation line to start with \code{\&} as well as ending
the previous line with one. fprettier treats a leading \code{\&} as a request
for \emph{manual alignment}: it keeps the column offsets the continuation lines
had relative to the first line, rather than aligning them itself. Whitespace
formatting still applies to the content of each line.

\begin{beforecol}
\begin{lstlisting}[style=before]
x = [ 1, 2, 3, &
    & 4, 5, 6 ]
\end{lstlisting}
\end{beforecol}
\begin{aftercol}
\begin{lstlisting}[style=after]
x = [1, 2, 3, &
    & 4, 5, 6]
\end{lstlisting}
\end{aftercol}

This is the right tool for tables of numbers, where the columns carry meaning.

\section{Turning whole features off}

\flag{no-indent} and \flag{no-whitespace} turn off the two main passes.
Case conversion, line splitting and comment spacing still run if enabled. In a
file that should never be reformatted, a directive does the same for just that
file:

\begin{lstlisting}[style=fortran]
! fprettier: --no-indent --no-whitespace
\end{lstlisting}

To skip a file altogether, exclude it with \flag{exclude}, or set
\flag{exclude-max-lines} to skip huge generated files.

\begin{exercise}[Which hatch?]
For each situation, pick the most appropriate mechanism: (a) a 30-line
\code{data} statement laid out as a matrix; (b) a vendored third-party module
you never want touched; (c) one \code{if} statement that you want on one line
even though it is 140 characters long; (d) a generated file of 200\,000 lines.
\end{exercise}
% ===========================================================================
\part{Inside fprettier}
% ===========================================================================

% ---------------------------------------------------------------------------
\chapter{Architecture}\label{ch:arch}
% ---------------------------------------------------------------------------

\section{The crate at a glance}

fprettier is a single Cargo package with a library target and a binary target.
The binary, \file{src/main.rs}, is a thin driver: it parses arguments, builds
configurations, finds files and runs the library's formatter on each one. All
the formatting logic lives in the library, whose public entry point is one
function:

\begin{lstlisting}[style=rust]
// src/process.rs
pub fn format_file<R: BufRead, W: Write>(input: R, output: &mut W, config: &Config) -> Result<()>
\end{lstlisting}

It reads Fortran from any buffered reader and writes the formatted text to any
writer. Table~\ref{tab:modules} maps the modules.

\begin{table}[ht]
\centering\small
\begin{tabularx}{\linewidth}{@{}l r X@{}}
\toprule
File & Lines & Responsibility\\
\midrule
\file{main.rs} & 650 & Driver: config per file, file discovery, parallelism, output modes, atomic writes\\
\file{cli.rs} & 730 & \code{clap} command definition and \code{CliArgs}\\
\file{config.rs} & 520 & \code{Config}, TOML discovery and merging, whitespace flag resolution\\
\file{directive.rs} & 260 & \code{! fprettier:} comment parsing\\
\file{process.rs} & 1710 & The pipeline: \code{format\_file}, \code{inspect\_file}, \code{format\_pass}\\
\file{scope.rs} & 270 & The \code{SCOPES} table of constructs that open a scope\\
\midrule
\file{parser/char\_filter.rs} & 380 & \code{CharFilter}: which characters are code, string, comment or fypp\\
\file{parser/stream.rs} & 560 & \code{InputStream}: physical lines to logical lines (\code{FortranLine})\\
\file{parser/patterns.rs} & 620 & About 75 compiled regular expressions for Fortran statements\\
\midrule
\file{format/indenter.rs} & 860 & \code{F90Indenter}: scope stack and indentation\\
\file{format/aligner.rs} & 400 & \code{F90Aligner}: continuation-line columns\\
\file{format/continuation.rs} & 230 & Manual alignment (leading \code{\&}) helpers\\
\file{format/whitespace.rs} & 1590 & The three-stage whitespace engine\\
\file{format/line\_split.rs} & 570 & Breaking long lines, detaching comments\\
\file{format/case\_convert.rs} & 1010 & Keyword tables and case conversion\\
\file{format/replacements.rs} & 170 & \code{.lt.} $\leftrightarrow$ \code{<}\\
\file{format/sort\_use.rs} & 460 & Sorting \code{use} statements and \code{only:} lists\\
\bottomrule
\end{tabularx}
\caption{Source files and what they do. Line counts include each file's unit
tests.}
\label{tab:modules}
\end{table}

The dependencies are deliberately few: \code{clap} for arguments,
\code{serde} and \code{toml} for configuration, \code{regex} for statement
recognition, \code{rayon} for parallelism, \code{walkdir} and \code{glob} for
finding files, \code{similar} for \flag{diff}, \code{anyhow} for errors, and
\code{mimalloc} as the global allocator.

\section{The flow of one file}

Figure~\ref{fig:pipeline} shows what happens to a file from the command line
to the disk. The rest of Part II takes the boxes one at a time.

\begin{figure}[ht]
\centering
\begin{tikzpicture}[node distance=5mm and 6mm]
  \node[stage] (cli) {\code{parse\_args}\\\file{cli.rs}};
  \node[stage,right=of cli] (collect) {\code{collect\_files}\\\file{main.rs}};
  \node[stage,right=of collect] (cfg) {\code{build\_config}\\per directory};
  \node[stage,below=10mm of cfg] (single) {\code{process\_single\_file}\\size, lines, directive};
  \node[stage,left=of single] (ff) {\code{format\_file}\\\file{process.rs}};
  \node[stage,left=of ff] (out) {write back /\\stdout / diff / check};

  \draw[flow] (cli) -- (collect);
  \draw[flow] (collect) -- (cfg);
  \draw[flow] (cfg) -- node[note,right]{rayon, one task per file} (single);
  \draw[flow] (single) -- (ff);
  \draw[flow] (ff) -- (out);

  % format_file internals
  \node[data,below=14mm of ff,minimum width=26mm] (buf) {bytes in memory};
  \node[stage,below=5mm of buf] (sort) {\code{sort\_use\_statements}\\(optional)};
  \node[stage,below=5mm of sort] (p1) {pass 1: \code{format\_pass}\\whitespace, no indent};
  \node[stage,right=12mm of p1] (insp) {\code{inspect\_file}\\requested indents};
  \node[stage,below=5mm of p1] (p2) {pass 2: \code{format\_pass}\\indent, no whitespace};
  \node[stage,below=5mm of p2] (crlf) {\code{write\_lines}\\restore CRLF};
  \draw[flow] (buf) -- (sort);
  \draw[flow] (sort) -- (p1);
  \draw[flow] (p1) -- (p2);
  \draw[flow] (p1.east) -- (insp.west);
  \draw[flow] (insp.south) |- (p2.east);
  \draw[flow] (p2) -- (crlf);
  \begin{scope}[on background layer]
    \node[draw=accent,dashed,rounded corners,fit=(buf)(sort)(p1)(insp)(p2)(crlf),inner sep=6pt,label={[note]north east:inside \code{format\_file}}] (box) {};
  \end{scope}
  \draw[densely dotted,thick,accent] (ff.south) -- (box.north -| ff.south);
\end{tikzpicture}
\caption{The life of a file. The lower box expands \code{format\_file}.}
\label{fig:pipeline}
\end{figure}

\section{Five ideas that explain the design}

\paragraph{1. Lines, not trees.} fprettier never builds a syntax tree. The
unit of work is the \emph{logical line}: one Fortran statement (or several
separated by semicolons), with its continuation lines and comments. Statements
are recognised by regular expressions anchored at the start of the logical
line; everything else is character-level work guided by bracket depth. This
keeps fprettier fast and makes it tolerant of code it does not fully
understand, at the price of a great many careful special cases.

\paragraph{2. A filter decides what is code.} Almost every component needs to
skip string literals, comments and fypp expressions. Rather than parse them
again and again, they all use one iterator adapter, \code{CharFilter}, which
yields only the characters of a line that are code (Chapter~\ref{ch:reading}).

\paragraph{3. Two passes.} Whitespace formatting and indentation run as two
separate passes over the whole file, the second reading the output of the
first. This separation keeps each pass simple, and it means the indentation pass
sees statements whose spacing is already final, which the continuation
aligner depends on (Chapter~\ref{ch:pipeline}).

\paragraph{4. A pre-pass reads the author's intent.} Before the indentation
pass, \code{inspect\_file} reads the existing indentation and records, for each
statement, how much indent its scope should get. This is how the stacking rule
of Section~\ref{sec:stacking} is implemented.

\paragraph{5. Output must be a fixed point.} Formatting fprettier's output
must change nothing. Any rule that measures the input (an original gap before a
comment, a manual alignment offset, a label's padding) must measure something
that its own output reproduces. Many of the comments in the source explain how
a piece of code keeps this property.

\section{Library use}

Because \code{format\_file} is generic over \code{BufRead} and \code{Write},
using fprettier from Rust takes a few lines. This program is the starting
point for many of the exercises in Part II:

\begin{lstlisting}[style=rust]
// examples/play.rs  (create it; run with `cargo run --example play`)
use fprettier::Config;
use fprettier::process::format_file;

fn main() {
    let src = "program p\ninteger::i\ni=1\nend program p\n";
    let config = Config { indent: 2, ..Config::default() };
    let mut out = Vec::new();
    format_file(src.as_bytes(), &mut out, &config).unwrap();
    print!("{}", String::from_utf8(out).unwrap());
}
\end{lstlisting}

A \code{\&[u8]} implements \code{BufRead}, so a byte string can be passed
directly. The crate also exports its components, such as \code{CharFilter}
from \code{fprettier::parser::char\_filter} and \code{F90Indenter} from
\code{fprettier::format::indenter}, which the integration tests drive one at a
time.

\begin{exercise}[Orientation]
Create \file{examples/play.rs} as above and run it. Then change it to read a
file named on its command line and format it with \code{indent: 3}. Compare the
result with \code{fprettier -s -i 3} on the same file.
\end{exercise}

\begin{exercise}[Reading the map]
Using Table~\ref{tab:modules} and \code{grep}, find the function that turns
\code{endif} into \code{end if}. Which file is it in, which configuration rule
switches it off, and which regular expression stops it from changing
\code{endif = 3}?
\end{exercise}

% ---------------------------------------------------------------------------
\chapter{Reading Fortran: CharFilter and InputStream}\label{ch:reading}
% ---------------------------------------------------------------------------

\section{The problem}

Consider this line:

\begin{lstlisting}[style=fortran]
write(*,*) 'Done! (100%)', n ! report ' progress
\end{lstlisting}

The first \code{!} is inside a string, so it does not start a comment; the
second does. The \code{(} inside the string is not a bracket, the \code{\%} is not a
component selector, and the quote inside the comment does not open a string. A
formatter that got any of these wrong would corrupt the program. Fortran's
lexical rules are simple enough that a small state machine gets them right,
and fprettier has exactly one: \code{CharFilter}.

\section{CharFilter}

\code{CharFilter} wraps a \code{CharIndices} iterator over a line and yields
\code{(byte\_position, char)} pairs, skipping characters according to three
switches:

\begin{lstlisting}[style=rust]
// src/parser/char_filter.rs
pub struct CharFilter<'a> {
    chars: std::iter::Peekable<std::str::CharIndices<'a>>,
    state: FilterState,
    filter_comments: bool,
    filter_strings: bool,
    filter_fypp: bool,
}

#[derive(Debug)]
struct FilterState {
    instring: StringDelimiter,
    infypp: bool,
    incomment: bool,
}

pub enum StringDelimiter {
    None,
    Single,     // '...'
    Double,     // "..."
    FyppHash,   // #{...}#
    FyppDollar, // ${...}$
    FyppAt,     // @{...}@
}
\end{lstlisting}

Two constructors cover nearly every use:

\begin{lstlisting}[style=rust]
/// Iterate the line's code: string literals, comments and fypp inline
/// blocks are all masked out
pub fn code(content: &'a str) -> Self { Self::new(content, true, true, true) }

/// Iterate the line's code and its comment, masking only string literals
/// and fypp inline blocks
pub fn code_and_comments(content: &'a str) -> Self { Self::new(content, false, true, true) }
\end{lstlisting}

The byte positions it yields are positions in the \emph{original} string, so a
caller can find a character with the filter and then slice the original line
around it. This is the pattern you will see everywhere in fprettier:

\begin{lstlisting}[style=rust]
/// Find the byte position of the first `!` that starts a comment,
/// i.e. outside strings and fypp expressions.
pub fn comment_start(line: &str) -> Option<usize> {
    CharFilter::code_and_comments(line).find_map(|(pos, c)| (c == '!').then_some(pos))
}
\end{lstlisting}

\subsection{The state machine}

Figure~\ref{fig:charfilter} shows the states. A \code{!} outside a string
enters the comment state, which lasts to the end of the line. A quote enters
the matching string state, and only the same quote leaves it. Fortran escapes
a quote inside a string by doubling it (\code{'it''s'}); the filter does not
special-case this, because closing and immediately reopening the string yields
the same classification of every character. A fypp opener
(\code{\#\{}, \code{\$\{}, \code{@\{}) enters a fypp state that only the matching
\code{\}\#}, \code{\}\$} or \code{\}@} leaves; quotes inside a fypp expression do not
open strings.

\begin{figure}[ht]
\centering
\begin{tikzpicture}[
  st/.style={rectangle,rounded corners=6pt,draw=accent,fill=accentlight,minimum width=24mm,minimum height=9mm,align=center,font=\sffamily\small},
  lab/.style={font=\sffamily\scriptsize,inner sep=2pt}]
  \node[st] (code) {code};
  \node[st,right=40mm of code] (dq) {in \texttt{"..."}};
  \node[st,above=9mm of dq] (sq) {in \texttt{'...'}};
  \node[st,below=9mm of dq] (fy) {in fypp \texttt{\$\{...\}\$}};
  \node[st,left=24mm of code] (cm) {comment};
  \draw[{Stealth[length=2.2mm]}-{Stealth[length=2.2mm]},thick,ink!70] (code.east) -- node[lab,above,sloped]{\texttt{\textquotesingle} \ldots{} \texttt{\textquotesingle}} (sq.west);
  \draw[{Stealth[length=2.2mm]}-{Stealth[length=2.2mm]},thick,ink!70] (code.east) -- node[lab,above]{\texttt{\textquotedbl} \ldots{} \texttt{\textquotedbl}} (dq.west);
  \draw[{Stealth[length=2.2mm]}-{Stealth[length=2.2mm]},thick,ink!70] (code.east) -- node[lab,below,sloped]{\texttt{\$\{} \ldots{} \texttt{\}\$} (also \texttt{\#}, \texttt{@})} (fy.west);
  \draw[flow] (code) -- node[lab,above]{!} (cm);
  \node[note,below=1mm of cm] {until end of line};
\end{tikzpicture}
\caption{CharFilter's states. A filtered state's characters (including the
delimiters) are skipped when the corresponding switch is on.}
\label{fig:charfilter}
\end{figure}

\subsection{Carrying state across lines}

A character literal may be continued onto the next line:

\begin{lstlisting}[style=fortran]
msg = 'this is a long message &
      &that continues here'
\end{lstlisting}

The second line begins \emph{inside} a string. \code{with\_string\_state}
creates a filter that starts in a given state, and \code{get\_string\_state}
reports the state at the end, so the stream reader can thread the state from
one physical line to the next.

\subsection{An iterative loop, on purpose}

The \code{next} method is a \code{loop} with \code{continue} for every skipped
character, not a recursive call. The comment in the source records why:

\begin{lstlisting}[style=rust]
fn next(&mut self) -> Option<Self::Item> {
    // Loops rather than recursing: a filtered run is as long as the string
    // or comment it covers, and one stack frame per skipped character
    // overflows the stack on a megabyte-long literal.
    loop {
        let (pos, c) = self.chars.next()?;
        // ...
\end{lstlisting}

\subsection{Splitting a line into masked regions}

Several rewriting passes (the context stage of the whitespace engine, case
conversion) need to transform the code of a line while leaving strings and
comments alone. \code{split\_masked\_regions} turns a line and a filter into a
list of alternating parts: even indices hold code, odd indices the masked text
verbatim, and concatenating the list reproduces the line exactly. A pass
rewrites the even parts and joins.

\begin{lstlisting}[style=plain]
line:    x = 'a, b' // y   ! note
parts:   ["x = ", "'a, b'", " // y   ", "! note"]     (with CharFilter::code)
\end{lstlisting}

\section{InputStream: from physical to logical lines}

\code{InputStream} reads a file one \emph{logical} line at a time. Its output
type is the single most important data structure in fprettier:

\begin{lstlisting}[style=rust]
// src/parser/stream.rs
pub struct FortranLine {
    /// The logical Fortran line (continuations joined, before semicolons split it)
    pub joined_line: String,
    /// Comments from each physical line
    pub comments: Vec<String>,
    /// The original physical lines
    pub lines: Vec<String>,
    /// OMP conditional prefix (e.g., "!$ ") that was stripped during parsing
    pub omp_prefix: String,
    /// For each physical line, whether it starts inside a string
    pub lines_in_string: Vec<bool>,
}
\end{lstlisting}

For this input

\begin{lstlisting}[style=fortran]
call compute(alpha, beta, &   ! first two
             gamma, delta)
\end{lstlisting}

\code{next\_fortran\_line} returns

\begin{lstlisting}[style=plain]
joined_line     = "call compute(alpha, beta, gamma, delta)"
lines           = ["call compute(alpha, beta, &   ", "             gamma, delta)"]
comments        = ["! first two", ""]
omp_prefix      = ""
lines_in_string = [false, false]
\end{lstlisting}

The \code{joined\_line} is what the scope tracker matches its regular
expressions against; the \code{lines} are what the formatter rewrites and
writes back. Note that \code{lines} keep their code only (the comment is in
\code{comments}), but keep their trailing blanks, which is how the gap before
an inline comment survives.

\subsection{Reading a physical line}

\code{read\_physical\_line} normalises one line of input:

\begin{enumerate}
  \item Tabs are expanded to eight spaces, but not inside a character
    context, where a tab is part of the constant's value (Fortran 2023,
    6.3.2.2). The string state from the previous physical line is passed in, so
    a tab in a continued literal is preserved too.
  \item The line ending, \code{\textbackslash n} or \code{\textbackslash r\textbackslash n}, is removed. (\code{format\_file}
    puts CRLF back at the very end; Section~\ref{sec:crlf}.)
  \item If the line starts inside a string and does not begin with \code{\&},
    an \code{\&} is prepended, so that the rest of the pipeline sees an
    ordinary continuation.
  \item An OpenMP conditional sentinel, \code{!\$} followed by a space, is
    split off into \code{omp\_prefix}. Without this, the sentinel would look
    like a comment and the code after it would be lost.
\end{enumerate}

\subsection{Assembling a logical line}

\code{next\_fortran\_line} then loops over physical lines:

\begin{enumerate}
  \item It splits the comment off with \code{comment\_start}.
  \item It runs a \code{CharFilter} over the code part to update the string
    state for the next line.
  \item A comment-only line in the middle of a continuation is kept (in
    \code{lines} and \code{comments}) but contributes no code, and reading
    continues.
  \item A fypp directive line in the middle of a continuation is treated the
    same way: preserved, but not part of the Fortran statement.
  \item A line that does not end with \code{\&} ends the logical line.
    Otherwise the trailing \code{\&} is removed and reading continues.
\end{enumerate}

Finally the code parts are joined with single spaces. Continuation parts are
trimmed, so the leading \code{\&} of a manually aligned line does not survive
into \code{joined\_line}.

\begin{pitfall}[Semicolons are not split]
\code{InputStream} does not split semicolon-separated statements, and a
comment in \code{next\_fortran\_line} explains why: splitting would turn one formatted
line into several. A logical line may therefore contain several statements,
and every consumer that matches statements (the indenter, the sort-use pass)
has to look for semicolons itself.
\end{pitfall}

\subsection{Guarding against pathological input}

Fortran 2023 removed the limit on the number of continuation lines, so
\code{InputStream} cannot reject a statement for having too many; the
standard bounds the statement's length instead, at one million characters.
\code{InputStream} allows ten times that, to leave room for indentation and
ampersands, and returns an error beyond it, so a corrupt or adversarial file
cannot exhaust memory.

\section{The pattern library}

\file{parser/patterns.rs} holds the regular expressions that recognise
statements. They are \code{static LazyLock<Regex>} values, compiled on first
use, built by a helper that makes them case-insensitive:

\begin{lstlisting}[style=rust]
fn build_re(pattern: &str) -> Regex {
    RegexBuilder::new(pattern)
        .case_insensitive(true)
        .unicode(true)
        .build()
        .unwrap_or_else(|_| panic!("invalid regex pattern: {pattern}"))
}

const EOL_STR: &str = r"\s*;?\s*$"; // End of line
const SOL_STR: &str = r"^\s*"; // Start of line

pub static IF_RE: LazyLock<Regex> = LazyLock::new(|| {
    build_re(&format!(r"{SOL_STR}(\w+\s*:)?\s*IF\s*\(.*\)\s*THEN(\s*$|\s*;)"))
});
pub static ENDIF_RE: LazyLock<Regex> =
    LazyLock::new(|| build_re(&format!(r"{SOL_STR}END\s*IF(\s+\w+)?{EOL_STR}")));
\end{lstlisting}

The patterns share a shape: start of line, an optional construct name
(\code{(\textbackslash w+\textbackslash s*:)?}), the keyword with
\code{\textbackslash s*} wherever Fortran allows the blank to be omitted
(\code{END\textbackslash s*IF} matches both \code{endif} and \code{end if}),
and an end anchor that allows a trailing semicolon. They are matched against
\code{CharFilter::code} output, so strings and comments cannot confuse them;
\code{IF\_RE}'s greedy \code{\textbackslash(.*\textbackslash)} is safe because
the line holds only code.

Some patterns are deliberately loose. \code{SUBR\_RE} accepts any prefix without
quotes before \code{SUBROUTINE}, which covers \code{recursive},
\code{pure}, \code{elemental}, \code{module} and their combinations without
listing them. The cost is that it also matches \code{end subroutine}, which
is why the indenter checks for \code{end} statements before it checks for
openers (Section~\ref{sec:indenter-order}).

\begin{keyidea}[Why pre-warm?]
There are about a hundred \code{LazyLock<Regex>} values across the crate. When
formatting in parallel, \file{main.rs} first formats a two-line program on the
main thread, so that all of them are compiled before the worker threads start
and race to initialise them. The comment there records the measured benefit
(10--40\,ms on a 3000-file run).
\end{keyidea}

\begin{exercise}[CharFilter by hand]
Write a unit test in \file{src/parser/char\_filter.rs} that collects
\code{CharFilter::code(line)} into a \code{String} for the line at the start of
this chapter. What does it contain? Now do the same with
\code{code\_and\_comments}.
\end{exercise}

\begin{exercise}[Doubled quotes]
Convince yourself that \code{'it''s'} is handled correctly by tracing the
state machine character by character. Then find the code in
\file{format/whitespace.rs} that handles doubled quotes explicitly. Why does
the whitespace engine need to, when \code{CharFilter} does not?
\end{exercise}

\begin{exercise}[Logical lines]
Use \code{InputStream::from\_string} in a test to print the
\code{FortranLine} for a statement with a comment-only line in the middle of
its continuation. Which vectors have three entries, and what does
\code{joined\_line} contain?
\end{exercise}
% ---------------------------------------------------------------------------
\chapter{The Formatting Pipeline}\label{ch:pipeline}
% ---------------------------------------------------------------------------

This chapter follows \code{format\_file} in \file{src/process.rs} from top to
bottom. It is the longest file in the crate because it is where every other
component is wired together, and where most of the special cases that tie them
together live.

\section{\texttt{format\_file}}

\begin{lstlisting}[style=rust]
pub fn format_file<R: BufRead, W: Write>(input: R, output: &mut W, config: &Config) -> Result<()> {
    let case_settings = CaseSettings::from_dict(&config.case_dict);
    let case_enabled = case_settings.is_enabled();
    let line_split_needed = config.line_length > 0 && config.line_length < LINE_SPLIT_THRESHOLD;
    let normalize_comment_spacing_needed = config.normalize_comment_spacing;

    let mut input_buffer = Vec::new();
    let mut reader = input;
    reader.read_to_end(&mut input_buffer)?;
    let crlf_input = input_buffer.windows(2).any(|pair| pair == b"\r\n");

    if input_buffer.len() > MAX_LINE_LENGTH {
        for warning in over_limit_warnings(&String::from_utf8_lossy(&input_buffer)) {
            eprintln!("Warning: {warning}");
        }
    }

    if config.sort_use || config.sort_use_only {
        if let Ok(src) = std::str::from_utf8(&input_buffer) {
            let sorted = sort_use_statements(src, config.sort_use, config.sort_use_only);
            input_buffer = sorted.into_bytes();
        }
    }
    // ... pass 1, inspection, pass 2 (below)
    write_lines(output, &formatted, crlf_input)?;
    Ok(())
}
\end{lstlisting}

The whole input is read into memory first, for three reasons: the CRLF check
needs to see it all, the \code{use} sorter rearranges whole lines, and the
inspection pre-pass reads the text a second time. (The driver has already read
the file into memory to look for a directive, so this costs a copy, not a
second read from disk.)

Notice that \code{sort\_use\_statements} runs \emph{first}, on raw text. A
collapsed \code{only:} list may come out longer than the line limit; running the
sorter first means the normal passes see the final line order and re-break
anything it left too long.

\subsection{The passes}

\begin{lstlisting}[style=rust]
    // Pass 1: Whitespace formatting
    let intermediate = if config.impose_whitespace {
        let mut pass1_output = Vec::new();
        format_pass(BufReader::new(Cursor::new(&input_buffer)), &mut pass1_output, config,
                    false, // Don't impose indent in pass 1
                    true,  // Do impose whitespace in pass 1
                    None)?;
        pass1_output
    } else {
        input_buffer
    };

    let inspect_result = if config.impose_indent && !config.strict_indent {
        Some(inspect_file(BufReader::new(Cursor::new(&intermediate)),
                          config.indent, config.strict_indent)?)
    } else {
        None
    };

    // Pass 2: Indentation
    let mut formatted = Vec::with_capacity(intermediate.len());
    if config.impose_indent {
        format_pass(BufReader::new(Cursor::new(intermediate)), &mut formatted, config,
                    true, false, inspect_result.as_ref())?;
    } else if !config.impose_whitespace
        && (case_enabled || line_split_needed || normalize_comment_spacing_needed)
    {
        format_pass(/* ... */ false, false, None)?;
    } else {
        formatted = intermediate;
    }
\end{lstlisting}

There is one function, \code{format\_pass}, that does both jobs, controlled by
two booleans. Pass 1 imposes whitespace, pass 2 imposes indentation. Features
that are not one or the other (case conversion, operator replacement, line
splitting and comment placement) run in \emph{every} pass that runs, which is
safe because each is idempotent. If both main passes are switched off but one
of those features is on, a single pass with both booleans false runs just
those features.

\begin{keyidea}[Inspect what pass 2 will read]
\code{inspect\_file} reads \code{intermediate}, pass 1's output, not the
original input. The comment in the source explains why: pass 1 collapses runs
of blank lines and can move a comment to its own line, so the original has a
different number of statements, and its requested indents would be paired with
the wrong lines. Pass 1 preserves each line's leading whitespace, which is all
the inspection reads.
\end{keyidea}

\subsection{Line endings}\label{sec:crlf}

\code{InputStream} strips \code{\textbackslash r} along with
\code{\textbackslash n}, so every formatted line ends with a bare
\code{\textbackslash n}. \code{write\_lines} puts \code{\textbackslash r\textbackslash n} back if the
input contained one anywhere. A file with mixed endings comes out entirely
CRLF.

\section{\texttt{inspect\_file}: reading the author's intent}\label{sec:inspect}

\code{inspect\_file} implements the stacking rule of
Section~\ref{sec:stacking}. It produces one number per non-blank logical line,
the \emph{requested indent}, plus the column of the first statement:

\begin{lstlisting}[style=rust]
struct InspectResult {
    /// Requested indents for each Fortran line
    required_indents: Vec<usize>,
    /// Indent level of first non-empty line
    first_indent: usize,
}
\end{lstlisting}

The heart of it:

\begin{lstlisting}[style=rust]
let offset = indent_of(first_line);
let mut required_indent = offset.saturating_sub(prev_offset);

if IF_RE.is_match(joined_trimmed) || DO_RE.is_match(joined_trimmed) {
    // The `indent_size > 0` guard is load-bearing: is_multiple_of(0)
    // is true only for 0 ...
    let indent_misaligned = indent_size > 0 && !offset.is_multiple_of(indent_size);
    if prev_offset != offset || strict_indent || indent_misaligned || prev_was_scope_opener {
        required_indent = indent_size;
    }
    // Otherwise keep required_indent as delta (usually 0 if same offset)
} else {
    required_indent = indent_size;
}

prev_was_scope_opener = opens_indented_scope(joined_trimmed);
required_indents.push(required_indent);
prev_offset = offset;
\end{lstlisting}

So every line requests a full indent, except an \code{if}/\code{do} line in
the same, aligned, column as the line before it (where the line before is not
one of the always-indented scope openers), which requests 0.

Blank lines are \emph{transparent}: they get no entry and do not update
\code{prev\_offset}. That answers Exercise~5.2: a blank line between
\code{program p} and \code{do} does not stop them from stacking.

\subsection{The off-by-one that is not a bug}

How is the requested indent used? In \code{format\_pass}, the line number used
to index \code{required\_indents} is incremented \emph{before} it is used, so
it starts at 1 and indexes a vector that starts at 0:

\begin{lstlisting}[style=rust]
// Count lines with content only, exactly as `inspect_file` records
// them. Note this leaves fortran_line_number 1-indexed against a
// 0-indexed vector, which shifts all values by 1: the alignment
// preservation logic depends on using the NEXT line's required_indent
// for scope-opening statements (IF/DO).
if !is_blank {
    fortran_line_number += 1;
}
\end{lstlisting}

and the value it reads is passed to the indenter as \code{relative\_indent},
the amount by which a \emph{scope opened on this line} indents its body. Put
together: the indent a scope gives its body is the requested indent of the
\emph{first statement of the body}. If that statement is a stacked
\code{do}, the scope's body indent is 0. That is the stacking rule.

\begin{figure}[ht]
\centering
\small
\begin{tabular}{@{}r l c c l@{}}
\toprule
$n$ & statement (column) & own entry & \code{relative\_indent} read & effect\\
\midrule
1 & \code{program p} (0)      & 4 & \code{[1]} = 0 & program body indent 0\\
2 & \code{do i=1,n} (0)       & 0 & \code{[2]} = 0 & loop body indent 0\\
3 & \code{if(\ldots)then} (0) & 0 & \code{[3]} = 4 & if body indent 4\\
4 & \code{b(i)=1} (0)         & 4 & \code{[4]} = 4 & (opens nothing)\\
5 & \code{end if} (0)         & 4 & \ldots & \\
\bottomrule
\end{tabular}
\caption{How the example of Section~\ref{sec:stacking} gets its
indentation. \emph{Own entry} is the line's \code{required\_indents[$n-1$]};
each line reads \code{required\_indents[$n$]}, the entry of the line after it.}
\label{fig:offbyone}
\end{figure}

\begin{pitfall}[Two loops that must agree]
\code{inspect\_file} and \code{format\_pass} each read the text with an
\code{InputStream} and each skip blank lines. If they ever disagreed about
which lines count, every requested indent after the first disagreement would be
paired with the wrong statement. Both therefore use the same predicate,
\code{is\_blank\_line}, and the comment on it says so.
\end{pitfall}

\section{\texttt{format\_pass}, step by step}

\code{format\_pass} creates an \code{InputStream} and, in the indentation pass,
an \code{F90Indenter}. Then, for each \code{FortranLine}, it runs the steps of
Figure~\ref{fig:format-pass}.

\begin{figure}[p]
\centering
\begin{tikzpicture}[node distance=3.2mm,
  st/.style={stage,minimum width=62mm,minimum height=6.5mm,font=\sffamily\scriptsize},
  side/.style={note,text width=52mm}]
  \node[st] (s1) {read \code{FortranLine}; skip repeated blank lines};
  \node[st,below=of s1] (s2) {copy \code{lines} $\to$ \code{output\_lines}; record \code{lines\_were\_indented}};
  \node[st,below=of s2] (s3) {\code{blank\_out\_omp\_prefix}};
  \node[st,below=of s3] (s4) {\code{split\_off\_label}};
  \node[st,below=of s4] (s5) {compute \code{FormattingFlags}};
  \node[st,below=of s5] (s6) {\code{extract\_and\_format\_pre\_ampersands}};
  \node[st,below=of s6] (s7) {\code{apply\_whitespace\_to\_lines}};
  \node[st,below=of s7] (s8) {\code{replace\_relational\_operators} (opt.)};
  \node[st,below=of s8] (s9) {\code{convert\_case} (opt.)};
  \node[st,below=of s9] (s10) {indentation: \code{compute\_and\_apply\_indentation}\\ or pin CPP line to column 0};
  \node[st,below=of s10] (s11) {\code{apply\_pre\_ampersand\_indentation}};
  \node[st,below=of s11] (s12) {\code{split\_lines\_if\_needed}};
  \node[st,below=of s12] (s13) {\code{compute\_comment\_indices}};
  \node[st,below=of s13] (s14) {\code{write\_output\_line} for each line};
  \foreach \a/\b in {s1/s2,s2/s3,s3/s4,s4/s5,s5/s6,s6/s7,s7/s8,s8/s9,s9/s10,s10/s11,s11/s12,s12/s13,s13/s14}{\draw[flow] (\a)--(\b);}
  \node[side,right=4mm of s3] {\code{!\$ } sentinel replaced by blanks};
  \node[side,right=4mm of s4] {numeric label removed, padding normalised (pass 2)};
  \node[side,right=4mm of s5] {fypp? CPP? \code{!\&} skip? auto-align?};
  \node[side,right=4mm of s6] {pass 2, manual alignment only};
  \node[side,right=4mm of s7] {pass 1 only; Chapter~\ref{ch:ws-engine}};
  \node[side,right=4mm of s10] {pass 2 only; Chapters \ref{ch:indenter}, \ref{ch:aligner}};
  \node[side,right=4mm of s12] {Section~\ref{sec:linesplit}};
  \node[side,right=4mm of s14] {label, comment placement, newline};
  \draw[flow] (s14.west) -- ++(-6mm,0) |- node[note,pos=0.25,left]{next line} (s1.west);
\end{tikzpicture}
\caption{What \code{format\_pass} does with each logical line.}
\label{fig:format-pass}
\end{figure}

\subsection{Steps 1--3: preparation}

Blank lines are dropped if the previous line was blank too, which is how runs of
blank lines collapse to one. (A labelled blank line does not count as blank
for this purpose: \code{skip\_blank = is\_blank \&\& labels.label.is\_empty()}.)

\code{lines\_were\_indented} records whether each physical line started with
whitespace \emph{before} anything rewrote it. It is needed at the very end, to
decide whether a comment-only line should be indented, and by then the lines
have been rewritten.

\code{blank\_out\_omp\_prefix} replaces the \code{!\$ } sentinel with three
spaces on every line, so that the formatting steps see plain Fortran in the
same columns. The indentation step writes the sentinel back
(\code{restore\_omp\_prefix} does it when indentation is off).

\subsection{Step 4: labels}

\code{split\_off\_label} removes a numeric statement label from the joined line and
from the first physical line, replacing it with blanks so that columns do not
move. It returns a \code{LineLabels}:

\begin{lstlisting}[style=rust]
struct LineLabels {
    joined_no_label: String,   // the joined logical line with label removed
    first_no_label: String,    // the first physical line with label removed
    first_label: String,       // the label extracted from first physical line
    label: String,             // the label extracted from joined line
    label_shift: usize,        // indentation shift due to label normalization
}
\end{lstlisting}

In the indentation pass it also normalises the run of blanks after the label,
and records in \code{label\_shift} how far that moved the statement, because
manually aligned continuation lines have to move by the same amount. It does
this only in pass 2, because pass 1's output is pass 2's input, and normalising
earlier would erase the original spacing that the shift is measured from. When
the line is finally written, \code{write\_output\_line} prints the label and
pads it to the statement's indentation column.

\subsection{Step 5: flags}

\begin{lstlisting}[style=rust]
let flags = FormattingFlags {
    is_fypp_line: is_fypp_directive(&fortran_line.joined_line),
    is_cpp_line: CPP_LINE_RE.is_match(&fortran_line.joined_line),
    skip_format: detect_skip_format(&fortran_line.comments, &mut in_deactivation_block),
    auto_align: should_auto_align(&output_lines),
};
\end{lstlisting}

\code{detect\_skip\_format} implements the \code{!\&} markers of
Chapter~\ref{ch:escape}, carrying the \code{!\&<} \ldots{} \code{!\&>} block
state in \code{in\_deactivation\_block} from line to line.
\code{should\_auto\_align} is false if any physical line starts with
\code{\&} followed by content: that is the request for manual alignment.
The helper \code{is\_rewritable\_code()} combines the first three flags: a
line may be rewritten only if it is neither a preprocessor line nor
deactivated.

\subsection{Steps 6--9: rewriting content}

Step 7, \code{apply\_whitespace\_to\_lines}, runs only in pass 1. For a
single-line statement it formats the label-free joined line. For a statement
with continuations it formats each physical line separately, removing the trailing
\code{\&} first and restoring it afterwards, and carries two pieces of state
from line to line: the bracket depth, so that \code{n=1} on the second line of
a call is still known to be a keyword argument, and the last significant
character, so that a \code{-} at the start of a continuation line can be
classified as unary or binary.

Manually aligned continuation lines cannot be formatted in pass 1, because
their leading \code{\&} is still attached. Step 6,
\code{extract\_and\_format\_pre\_ampersands}, runs in pass 2: it records each
line's manual offset, strips the leading \code{\&} with
\code{remove\_pre\_ampersands}, and then applies whitespace formatting to the
continuation lines. Chapter~\ref{ch:aligner} covers this.

Steps 8 and 9 apply operator replacement and case conversion to each physical
line, through \code{map\_code\_lines}, which skips lines that begin inside a
multi-line string.

\subsection{Steps 10--11: indentation}

In pass 2, a C preprocessor line is pinned to column 0 and the indenter is
\emph{not} called, since a directive does not affect Fortran scope. Every other line
goes to \code{compute\_and\_apply\_indentation}, which calls
\code{F90Indenter::process\_logical\_line} and applies the resulting
per-line indents. It is full of special cases, each documented in place:

\begin{itemize}
  \item a labelled first line keeps its leading blanks, which encode the
    label's column for \code{write\_output\_line};
  \item continuation lines of a multi-line fypp directive keep their offset
    relative to the directive's first line (\code{rebase\_fypp\_continuation});
    using their absolute indent instead made such blocks drift right on every
    run, a bug fixed in 0.4.0;
  \item a continuation line starting with \code{\&} is indented one column less, so
    that the content after the \code{\&} lands on the alignment column;
  \item an OpenMP conditional line gets its sentinel back at column 0, with
    padding after it.
\end{itemize}

Step 11 handles manually aligned lines, putting their leading \code{\&} back with
\code{prepend\_ampersands} and indenting them.

\subsection{Steps 12--14: splitting and writing}

\code{split\_lines\_if\_needed} breaks over-long lines (never preprocessor
lines, whose continuation syntax differs), returning for each output line the
index of the input line it came from. \code{compute\_comment\_indices} uses that
mapping to decide which output line carries each comment: the last fragment of
its origin line. Finally \code{write\_output\_line} writes each line with its
label and comment. Chapter~\ref{ch:lines} covers both.

\section{Context structs}

\code{format\_pass} passes a lot of state to its helpers. Rather than ten-argument
functions, it bundles it:

\begin{description}[font=\ttfamily\small]
  \item[PassContext] per pass: the config, the resolved
    \code{whitespace\_flags: [bool; 11]}, which pass this is, and the
    \code{InspectResult};
  \item[FormattingFlags] per logical line: the four flags above, and
    \code{Copy};
  \item[LineLabels] per logical line: the label strings;
  \item[LineWriteContext] per logical line, for writing: computed indents,
    \code{lines\_were\_indented}, comment indices, split origins and the line
    length.
\end{description}

\begin{exercise}[Watch the passes]
In a scratch copy of \file{process.rs}, make \code{format\_file} print
\code{intermediate} to standard error between the passes. Format
\file{examples/in/example.f90}. Which changes happened in pass 1, and which in
pass 2? Find a line that pass 2 changed in a way you did not expect.
\end{exercise}

\begin{exercise}[Requested indents]
Add an \code{eprintln!} of \code{inspect\_result} and format the stacking example
of Section~\ref{sec:stacking}. Check your output against
Figure~\ref{fig:offbyone}.
\end{exercise}

\begin{exercise}[Inspect the wrong buffer]
Change \code{format\_file} so that \code{inspect\_file} reads the original input
instead of \code{intermediate} (you will need to clone \code{input\_buffer}
before pass 1 consumes it), and run \code{cargo test}. Exactly one test
fails: \code{test\_stacking\_survives\_a\_detached\_comment\_above} in
\file{tests/integration\_test.rs}. Before writing this book, no test failed;
that test was added when the gap was found. Read it, and explain line by line
how the requested indents end up paired with the wrong statements. (Hint: pass
1 turns a long inline comment into a line of its own, which adds a logical
line.)
\end{exercise}
% ---------------------------------------------------------------------------
\chapter{Scopes and the Indenter}\label{ch:indenter}
% ---------------------------------------------------------------------------

\section{The scope table}

Every construct that opens an indented scope is one row of a static table in
\file{src/scope.rs}:

\begin{lstlisting}[style=rust]
pub struct Scope {
    pub group: Group,
    /// Statement that opens the scope
    pub open: Option<&'static Re>,
    /// Statement that continues it at the parent indent (ELSE, CASE, CONTAINS)
    pub cont: Option<&'static Re>,
    /// Statement that closes it
    pub close: Option<&'static Re>,
    /// Whether a close must match this scope exactly; a bare END matches any
    pub spec: bool,
    /// WHERE/FORALL, which open a block only when nothing follows the `)`
    pub conditional: bool,
}

pub static SCOPES: &[Scope] = {
    use Group::{Core, Fypp, Module};
    &[
        // 0: IF/THEN/ELSE/ENDIF
        Scope::new(Core, Some(&IF_RE), Some(&ELSE_RE), Some(&ENDIF_RE)),
        // 1: DO/ENDDO
        Scope::new(Core, Some(&DO_RE), None, Some(&ENDDO_RE)),
        // 2: SELECT CASE/RANK/TYPE
        Scope::new(Core, Some(&SELCASE_RE), Some(&CASE_RE), Some(&ENDSEL_RE)),
        // ... 25 rows in all; Appendix C lists them
    ]
};
\end{lstlisting}

A scope's identity is its \emph{index} in this table, which is what the
indenter pushes on its stack, so rows may be appended but never reordered.
Row 9 is special: it has no opener and its closer is \code{ENDANY\_RE}, a bare
\code{end}, with \code{spec: false}, meaning that it may close any scope.

Each row belongs to a \code{Group}. \code{Core} rows are always active;
\code{Module} rows (\code{module}, \code{submodule}, \code{program}) only with
\code{indent\_mod}; \code{Fypp} rows only with \code{indent\_fypp} in the
indentation pass. A \code{ScopeParser} is nothing more than those two
booleans, with an iterator over the enabled rows:

\begin{lstlisting}[style=rust]
impl ScopeParser {
    fn enabled(self, idx: usize) -> bool {
        match SCOPES.get(idx).map(|s| s.group) {
            Some(Group::Core) => true,
            Some(Group::Module) => self.indent_mod,
            Some(Group::Fypp) => self.indent_fypp,
            None => false,
        }
    }
    pub fn iter(self) -> impl Iterator<Item = (usize, &'static Scope)> {
        SCOPES.iter().enumerate().filter(move |(idx, _)| self.enabled(*idx))
    }
}
\end{lstlisting}

The table is the single source of truth for which constructs open scopes.
Before it existed, \code{inspect\_file} kept its own hand-written copy of the
openers, and missed some; \code{opens\_indented\_scope} now derives its list
from the same patterns.

\section{\texttt{F90Indenter}'s state}

\begin{lstlisting}[style=rust]
pub struct F90Indenter {
    /// Stack of open scopes, as indices into SCOPES
    scope_storage: Vec<usize>,
    /// Stack of absolute indent levels
    indent_storage: Vec<usize>,
    /// Output: indent for each physical line in current logical_line
    line_indents: Vec<usize>,
    parser: ScopeParser,
    aligner: F90Aligner,
    initial: bool,
    /// Open labeled DO constructs, as (scope_storage depth after the push,
    /// terminating statement label)
    do_labels: Vec<(usize, String)>,
}
\end{lstlisting}

The two stacks move together. \code{indent\_storage} starts with one entry,
the column of the file's first statement (\code{first\_indent} from
\code{inspect\_file}), so it always has one entry more than
\code{scope\_storage}. Its top is the indent of the current scope's body;
the entry below it is the indent of the statement that opened the scope.

Table~\ref{tab:stacks} traces both stacks through the \file{geometry.f90}
example of Chapter~\ref{ch:start}.

\begin{table}[ht]
\centering\small
\begin{tabular}{@{}l c l l@{}}
\toprule
statement & indent & \code{scope\_storage} after & \code{indent\_storage} after\\
\midrule
\code{module geometry}  & 0  & [13]          & [0, 4]\\
\code{implicit none}    & 4  & [13]          & [0, 4]\\
\code{contains}         & 0  & [13]          & [0, 4]\\
\code{subroutine area}  & 4  & [13, 3]       & [0, 4, 8]\\
\code{real, \ldots}     & 8  & [13, 3]       & [0, 4, 8]\\
\code{if (r < 0.0) then}& 8  & [13, 3, 0]    & [0, 4, 8, 12]\\
\code{a = -1.0}         & 12 & [13, 3, 0]    & [0, 4, 8, 12]\\
\code{else}             & 8  & [13, 3, 0]    & [0, 4, 8, 12]\\
\code{a = pi*r**2}      & 12 & [13, 3, 0]    & [0, 4, 8, 12]\\
\code{end if}           & 8  & [13, 3]       & [0, 4, 8]\\
\code{end subroutine}   & 4  & [13]          & [0, 4]\\
\code{end module}       & 0  & []            & [0]\\
\bottomrule
\end{tabular}
\caption{The indenter's stacks, line by line. Scope 13 is \code{MODULE}, 3 is
\code{SUBROUTINE}, 0 is \code{IF}.}
\label{tab:stacks}
\end{table}

Three rules produce the \emph{indent} column:

\begin{itemize}
  \item an ordinary statement, or one that opens a scope, is written at the top of
    \code{indent\_storage};
  \item a statement that continues a scope (\code{contains}, \code{else})
    is written one entry down, at the indent of the scope's opener;
  \item a valid \code{end} is written one entry down too, and then the scope is
    popped.
\end{itemize}

\section{\texttt{process\_logical\_line}}

This method is called once per logical line. Its outline:

\begin{lstlisting}[style=rust]
pub fn process_logical_line(&mut self, logical_line: &str, lines: &[String],
                            params: &IndentParams<'_>) {
    self.line_indents.clear();
    let filtered_line: String = CharFilter::code(logical_line).map(|(_, c)| c).collect();
    let parts: Vec<&str> = filtered_line.split(';').collect();

    let end = self.detect_end_statements(&parts, &filtered_line, params.indent_fypp);
    let (is_continue, fypp_continue_pop_to) = self.detect_continuation(&filtered_line);
    let (new_scope, additional_scopes) = if end.is_any || is_continue {
        (None, Vec::new())
    } else {
        self.detect_new_scopes(&parts, &filtered_line)
    };

    let line_indent = self.calculate_indent(is_continue, end.is_any,
                                            end.valid || end.valid_after_semicolon);
    // ... continuation lines: aligner or manual offsets, added to a base indent
    // ... pop indent_storage for ENDs, push for new scopes
    self.close_labeled_do(&filtered_line, params.label, new_scope.is_some());
    self.initial = false;
}
\end{lstlisting}

The statement is matched with strings and comments removed, and split on
semicolons, so that each statement of \code{x = 1; if (y) then} is examined.

\subsection{Order matters}\label{sec:indenter-order}

The three detections run in a fixed order: \code{end}, then continuation, then
new scope, and a line that matched an earlier one is not tested for a new scope.
Two comments in the source give the reasons. First, \code{SUBR\_RE} and
\code{FCT\_RE} match \code{end subroutine} and \code{end function}, because
they accept any prefix. Second, some continuations look like openers:
\code{type is (t)} matches \code{TYPE\_RE} as a derived-type definition with a
type-parameter list. Since no Fortran statement both continues a construct and
opens a new one, testing for a new scope last, and only if nothing else matched,
avoids both traps.

\subsection{Detecting ends}

\code{detect\_end\_statements} tries every enabled closer against every
semicolon-separated part. For the first part, it pops \code{scope\_storage}
and checks the popped scope against the closer:

\begin{lstlisting}[style=rust]
let popped_scope = self.scope_storage.pop().expect("...");
let popped_fypp_scope = SCOPES[popped_scope].is_fypp();
if !scope.spec || popped_scope == scope_idx || (indent_fypp && popped_fypp_scope) {
    end.valid = true;
} else {
    // The END closes a construct we never opened (an opener pattern we don't
    // recognize). Put the scope back: dropping it would desync
    // scope_storage from indent_storage and dedent the rest of the file.
    self.scope_storage.push(popped_scope);
}
\end{lstlisting}

An \code{end} is valid if it is a bare \code{end} (\code{spec} false), if it
matches the innermost scope, or if the innermost scope is a fypp block (so that
Fortran code closing a construct opened in a fypp branch does not wreck the
stacks). An \code{end} that matches nothing is not valid, and the scope goes
back. This is the fix described in the 0.5.0 changelog: before it, an
unrecognised \code{end} popped anyway, and everything after it was
under-indented.

For later parts (after a semicolon), an \code{end} is counted in
\code{additional\_count}, unless the same line also contains its opener
(\code{do i=1,n; x=x+i; end do} is balanced and changes nothing).

\subsection{Detecting continuations}

A Fortran continuation (\code{else}, \code{case}, \code{contains},
\code{elsewhere}) counts only if the scope it continues is on \emph{top} of the
stack: \code{else} inside a \code{do} inside an \code{if} is not the
\code{if}'s \code{else}. A fypp continuation (\code{\#:else}, \code{\#:elif}) is
different: fypp directives may enclose Fortran constructs, so
\code{detect\_continuation} searches the whole stack for the matching fypp
scope, and returns the depth to pop to after the line's indent is computed.
Only \code{scope\_storage} is popped, not \code{indent\_storage}, so the body
of the \code{\#:else} branch is indented like the body of the \code{\#:if}
branch.

\subsection{Detecting new scopes}

\code{detect\_new\_scopes} takes the first enabled row whose opener matches
(table order breaks ties), with one twist for \code{where} and \code{forall}:

\begin{lstlisting}[style=rust]
if open.is_match(filtered_line) {
    // WHERE/FORALL only open a scope if nothing follows the
    // closing parenthesis: "WHERE (x > 0)" opens one, but
    // "WHERE (x > 0) y = 1" is a single statement.
    if !scope.conditional || Self::is_where_forall_block(filtered_line) {
        new_scope = Some(scope_idx);
    }
    break;
}
\end{lstlisting}

\code{is\_where\_forall\_block} walks the line with bracket counting to find the
parenthesis that closes the mask and checks what follows it. It is the one
caller that builds a \code{CharFilter} with fypp filtering \emph{off}, because a
\code{\$\{\ldots\}\$} after the parenthesis is content, which makes the statement
a one-liner.

Openers after semicolons are collected separately into
\code{additional\_scopes}, each pushed with a full indent.

\subsection{Pushing}

A new scope is pushed with the line's own indent plus \code{relative\_indent}:

\begin{lstlisting}[style=rust]
fn push_scope(&mut self, scope: usize, current_indent: usize, relative_indent: usize) {
    self.scope_storage.push(scope);
    self.indent_storage.push(current_indent + relative_indent);
}
\end{lstlisting}

\code{relative\_indent} is the requested indent from \code{inspect\_file}
(Section~\ref{sec:inspect}), which is normally the indent width and 0 for a
stacked scope.

\subsection{Labelled DO loops}

\code{do 100 i = 1, 10} has no \code{end do}; it ends at the statement labelled
100. When such a \code{do} opens a scope, \code{close\_labeled\_do} records the
label and the stack depth in \code{do\_labels}. When a later line carries a
label, every recorded loop with that label is closed, innermost first, which
handles several loops sharing one terminator. Because the terminating statement
belongs to the loop body, this happens \emph{after} the line's own indent has
been computed.

\section{Continuation lines}

When a logical line has more than one physical line, the indenter computes one
indent per physical line. The first line gets \code{line\_indent}. The rest get
an offset relative to the start of the statement, from either the manual
offsets (lines with a leading \code{\&}, Section~\ref{sec:manual}) or the
aligner (Chapter~\ref{ch:aligner}), plus a base:

\begin{lstlisting}[style=rust]
let base_indent = if (end.is_any && end.valid) || is_continue {
    // For valid END and CONTINUE, use grandparent indent
    if self.indent_storage.len() >= 2 {
        self.indent_storage[self.indent_storage.len() - 2]
    } else {
        *self.indent_storage.last().unwrap_or(&0)
    }
} else if new_scope.is_some() && params.manual_lines_indent.is_some() {
    line_indent
} else {
    *self.indent_storage.last().unwrap_or(&0)
};
for &align_offset in continuation_indents.iter().skip(1) {
    self.line_indents.push(align_offset + base_indent);
}
\end{lstlisting}

The base is computed before the stacks are updated, so a continued
\code{end} or \code{else} line aligns its continuations relative to its own
column.

\begin{exercise}[Trace the stacks]
Write a unit test in \file{src/format/indenter.rs} that feeds the statements of
the \code{select case} example from Chapter~\ref{ch:tour} to an indenter, as
\code{test\_simple\_if} does, and asserts the indents. Add a temporary
\code{eprintln!} of both stacks at the end of \code{process\_logical\_line}
and compare with Table~\ref{tab:stacks}.
\end{exercise}

\begin{exercise}[An end that closes nothing]
Predict the indentation of this fragment inside a subroutine, where an
\code{end associate} has lost its \code{associate}:
\begin{lstlisting}[style=fortran]
do i = 1, n
x = 1
end associate
y = 2
end do
z = 3
\end{lstlisting}
Check your prediction. Which statement in \code{detect\_end\_statements} keeps
\code{y = 2} inside the loop and \code{z = 3} at the subroutine's level? What
would happen to the rest of the file without it?
\end{exercise}

\begin{exercise}[A new scope]
Add a scope for a hypothetical \code{begin foo} \ldots{} \code{end foo} construct:
two patterns in \file{parser/patterns.rs}, one row appended to \code{SCOPES},
and a test. Should it also go into \code{opens\_indented\_scope}? (See
Chapter~\ref{ch:recipes} for the full recipe.)
\end{exercise}

% ---------------------------------------------------------------------------
\chapter{Aligning Continuation Lines}\label{ch:aligner}
% ---------------------------------------------------------------------------

\section{The aligner's model}

\code{F90Aligner} answers one question: for each physical line after the first,
at which column, relative to the start of the statement, should it start? It
walks the physical lines of a logical line, character by character through
\code{CharFilter::code}, maintaining a stack of candidate columns:

\begin{lstlisting}[style=rust]
pub struct F90Aligner {
    line_indents: Vec<usize>,
    level: usize,                     // Bracket nesting level
    bracket_indent_stack: Vec<usize>, // Stack of alignment positions
    /// Whether an assignment in this logical line has set the alignment
    assign_aligned: bool,
}
\end{lstlisting}

At the end of each physical line, the top of \code{bracket\_indent\_stack}
becomes the indent of the next line. The stack is pushed and popped by these
events:

\begin{description}[font=\sffamily\small\bfseries]
  \item[Opening bracket] \code{(}, \code{(/} or \code{[} at position $p$
    pushes $p + \mathrm{len} + r$, the column just inside it, where $r$ is the
    current top (so positions on continuation lines are made relative to the
    statement start).
  \item[Closing bracket] pops.
  \item[Assignment] An \code{=} (or \code{=>}) at depth zero, not part of a
    relational operator, pushes the column after it and its following space.
    Only the first one in a statement counts (\code{assign\_aligned}); before
    0.5.0, every continuation line's \code{=} pushed another column and each
    added to the last, so the indentation ran away.
  \item[Top-level comma] after an assignment pops the assignment's column
    again, so the second item of \code{a = 1, b = 2} is not aligned to the
    first \code{=}.
  \item[Declaration] In a declaration (detected with \code{VAR\_DECL\_RE},
    \code{PUBLIC\_RE} or \code{PRIVATE\_RE}), \code{::} pushes the column after
    \code{:: }.
  \item[\texttt{use}] In a \code{use} statement, a single colon (the one after
    \code{only}) pushes the column after \code{: }.
\end{description}

Two rules then adjust the result at the end of each line. If the line ends with an
opening bracket followed by \code{\&} (\code{call foo( \&}), the bracket's
column is replaced by the enclosing one, since aligning to a column at the very
end of the line would waste it; the same goes for an \code{=}, \code{::} or
\code{:} directly before the break, which are simply not pushed. And if the
top of the stack is 0, meaning there is nothing to align to, it is replaced by
the indent width.

\section{Worked examples}

The alignment columns for the examples of Chapter~\ref{ch:tour}, inside a
program whose body is at column 4:

\begin{center}
\small
\begin{tabular}{@{}l l r r@{}}
\toprule
first line & event & relative & absolute\\
\midrule
\code{call compute(alpha, beta, \&} & \code{(} at 12 & 13 & 17\\
\code{total = first\_value + \&} & \code{=} at 6 & 8 & 12\\
\code{integer, dimension(3) :: a, \&} & \code{::} at 22 & 25 & 29\\
\code{x = [1, 2, \&} & \code{[} at 4 & 5 & 9\\
\code{call foo( \&} & \code{(} before break & 4 (default) & 8\\
\bottomrule
\end{tabular}
\end{center}

In the \code{x = [1, 2, \&} case the assignment pushes 4 first, and the bracket
then pushes 5 on top of it; the innermost open bracket wins.

\begin{keyidea}[Measure the final text]
The aligner measures columns in the lines it is given, so it must be given
lines whose spacing is final: had it seen \code{call  compute( alpha,}, it would
put the continuation two columns too far right, and the whitespace engine would
then shrink the first line without moving the continuation. The ordering
guarantees it: the indentation pass reads pass 1's output, and even within one
\code{format\_pass}, whitespace (step 7 of Figure~\ref{fig:format-pass}) runs
before indentation (step 10).
\end{keyidea}

\section{Manual alignment}\label{sec:manual}

When any continuation line starts with \code{\&} followed by content,
\code{should\_auto\_align} returns false and the aligner is not used.
Instead, \file{format/continuation.rs} provides three functions, used in
pass 2 by \code{process.rs}:

\begin{description}[font=\ttfamily\small]
  \item[get\_manual\_alignment] computes each line's indent (blanks plus
    \code{\&}) relative to the first line. These offsets replace the aligner's.
    A special case normalises ``simple'' continuations, where the first
    continuation line has no leading \code{\&}: their offsets are shifted so the
    first one equals the continuation indent.
  \item[remove\_pre\_ampersands] strips the leading \code{\&} and blanks from
    each line, remembering the stripped prefix (\code{"\& "}, \code{"\&   "})
    and how many blanks the previous line had before its trailing \code{\&}.
    Lines that are fypp directives or begin inside a string are left alone,
    since a leading \code{\&} there is content.
  \item[prepend\_ampersands] puts the prefixes back after indentation,
    reducing the computed indent by one for the \code{\&} itself.
\end{description}

Between removing and putting back the prefixes, pass 2 runs the whitespace engine
over the continuation lines, which pass 1 had to skip. It reproduces the
bracket-depth and last-character tracking of pass 1, and restores each trailing
\code{\&} with the spacing recorded in \code{ampersand\_sep}.

\begin{exercise}[Columns by hand]
Compute, by hand, the continuation column fprettier will choose for
\begin{lstlisting}[style=fortran]
    ptr => make_node(left=build(a, &
    b), right=c)
\end{lstlisting}
then check. Which pushes and pops happen on the first line?
\end{exercise}

\begin{exercise}[Comma resets]
Find the code that pops the assignment column on a top-level comma, and
comment it out, and format \code{a = 1, \&} / \code{b = 2}. What changes?
Then think of a real Fortran statement with a top-level \code{=} followed by a
top-level comma (hint: \code{do}). Is the rule right for it?
\end{exercise}
% ---------------------------------------------------------------------------
\chapter{The Whitespace Engine}\label{ch:ws-engine}
% ---------------------------------------------------------------------------

\section{Three stages}\label{sec:ws-stages}

\file{src/format/whitespace.rs} formats one line in three stages:

\begin{lstlisting}[style=rust]
pub fn format_line_with_level(
    logical_line: &str,
    whitespace_flags: &[bool; 11],
    format_decl: bool,
    initial_level: usize,
    prev_line_last_char: Option<char>,
) -> (String, usize, Option<char>) {
    let mut line = logical_line.to_string();
    // Stage 1: Remove extra whitespace
    line = rm_extra_whitespace(&line, format_decl);
    // Stage 2: Add whitespace character-wise (with bracket level tracking)
    let (formatted, final_level, last_char) = add_whitespace_charwise_with_level(
        &line, whitespace_flags, format_decl, initial_level, prev_line_last_char);
    line = formatted;
    // Stage 3: Add whitespace context-aware
    line = add_whitespace_context(&line, whitespace_flags);
    (line, final_level, last_char)
}
\end{lstlisting}

The extra parameters and results exist for continuation lines: the bracket
depth at the start of the line and at its end, and the last significant
character of the previous line and of this one.

Figure~\ref{fig:ws-stages} follows one line through the stages at the default
level.

\begin{figure}[ht]
\centering
\begin{tikzpicture}[node distance=4mm]
  \node[data,minimum width=70mm] (in) {x\ \ =\ \ a( i , j )+b\ \ *c};
  \node[data,minimum width=70mm,below=7mm of in] (s1) {x=a(i,j)+b*c};
  \node[data,minimum width=70mm,below=7mm of s1] (s2) {x = a(i, j) + b*c};
  \node[data,minimum width=70mm,below=7mm of s2] (s3) {x = a(i, j) + b*c};
  \draw[flow] (in) -- node[note,right=2mm]{stage 1: \code{rm\_extra\_whitespace}} (s1);
  \draw[flow] (s1) -- node[note,right=2mm]{stage 2: \code{Charwise} handlers} (s2);
  \draw[flow] (s2) -- node[note,right=2mm]{stage 3: regex-driven context rules} (s3);
\end{tikzpicture}
\caption{One line through the three stages at level 2.}
\label{fig:ws-stages}
\end{figure}

\subsection{Stage 1: remove}

\code{rm\_extra\_whitespace} walks the code characters (via \code{CharFilter},
copying skipped strings and comments through verbatim) and keeps a blank only
where it separates two words: after an alphanumeric character, and not before a
non-word character. It preserves leading whitespace, which the aligner and the
inspection pass need. It also keeps blanks in a few places where removing them
would change a token's meaning or a deliberate layout:

\begin{itemize}
  \item around \code{::} when \code{format\_decl} is off, so hand-aligned
    declaration columns survive;
  \item before a string literal (\code{print *, 'x'});
  \item next to fypp inline expressions, whose content is opaque.
\end{itemize}

\subsection{Stage 2: the \texttt{Charwise} handlers}

Stage 2 is a small hand-written scanner. Its state is a struct, and its main
loop offers each character to a sequence of handlers; the first that recognises
it writes output, advances the cursor, and returns true:

\begin{lstlisting}[style=rust]
while f.i < f.chars.len() {
    let handled = f.try_verbatim_regions()
        || f.try_open_delimiter()
        || f.try_close_delimiter()
        || f.try_comma_semicolon()
        || f.try_type_component()
        || f.try_assignment()
        || f.try_declaration()
        || f.try_concat()
        || f.try_plus_minus()
        || f.try_mult_div();
    if !handled {
        f.out.push(f.chars[f.i]);
        f.i += 1;
    }
}
\end{lstlisting}

The order matters. \code{try\_verbatim\_regions} must come first, so that
nothing inside a string or a fypp expression is reformatted.
\code{try\_concat} must come before \code{try\_mult\_div}, or \code{//} would be
read as two divisions. Each handler follows the same pattern: it removes
trailing blanks from the output with \code{pop\_trailing\_spaces}, writes the
token with the spacing it wants, and skips any blanks in the input after it.
A handler that does not apply returns false without touching the cursor.

Several handlers depend on the bracket \code{level}: \code{try\_assignment}
spaces \code{=} only at level 0 (pointer assignment at any level), and
\code{needs\_space\_before\_delimiter} uses it to put a space after a
top-level assignment. The level starts at \code{initial\_level}, which is
how a keyword argument on a continuation line is still recognised as being
inside brackets.

\paragraph{Unary or binary?} \code{try\_plus\_minus} has the most interesting
decision:

\begin{lstlisting}[style=rust]
let prev_char = self.prev_non_space();
let is_exponent = match prev_char {
    Some(c) if c == 'd' || c == 'D' || c == 'e' || c == 'E' => {
        let trimmed = self.out.trim_end();
        SCI_NOTATION_RE.is_match(trimmed)
    }
    _ => false,
};
let is_binary = match prev_char {
    Some(c) => (c.is_alphanumeric() || c == ')' || c == ']') && !is_exponent,
    None => {
        if self.out.trim().is_empty() {
            match self.prev_line_last_char {
                Some(c) => c.is_alphanumeric() || c == ')' || c == ']',
                None => false,
            }
        } else {
            false
        }
    }
};
\end{lstlisting}

A sign is binary when it follows an operand. The exponent check uses
\code{SCI\_NOTATION\_RE}, \code{(?:\^{}|[\^{}a-z\_])(\textbackslash
d+\textbackslash.?\textbackslash d*|\textbackslash d*\textbackslash.?\textbackslash
d+)[ed]\$}, which requires the \code{e} or \code{d} to end a number that is not
the tail of a name: \code{1.0d-3} is an exponent, \code{val\_1d-2} is not.

\subsection{Stage 3: context rules}

\code{add\_whitespace\_context} splits the line into code and masked parts with
\code{split\_masked\_regions} and applies regular-expression rewrites to the
code parts: relational operators, then the \code{?}/\code{:} of conditional
expressions (with a per-bracket stack of ``seen a \code{?}'' flags, carried
across masked parts), logical operators, user-defined operators (skipping
\code{.true.} and \code{.false.}), binary plus and minus, and \code{print *}.
A pointer assignment's \code{ => } is swapped for a placeholder during the
relational rewrite, so that its \code{>} is not mistaken for greater-than.

Finally \code{space\_out\_keywords} applies rules that span code and string parts,
so run on the rejoined line: \code{namelist} groups, \code{endif} $\to$
\code{end if} (guarded by \code{END\_WITH\_EQ\_RE} against \code{endif = 3}),
the same after a semicolon, \code{only: }, and construct names.

\section{A quirk, and its fix}

Until November 2026, the line of Figure~\ref{fig:ws-stages} came out at level 0
as \code{x=a(i,j) +b*c}, with a blank on one side of the \code{+} only. The
cause shows how the handlers depend on one another. Stage 1 produces
\code{x=a(i,j)+b*c}. In stage 2 the closing parenthesis goes to
\code{try\_close\_delimiter}, which writes \code{)} and then decides what
spacing follows. Its list of characters that take no blank did not include
\code{+}, so it wrote one. At level 2, \code{try\_plus\_minus} then saw the
\code{+}, popped that trailing blank and wrote \code{ + } symmetrically. At level
0, or with \code{-{}-whitespace-plusminus=false}, that handler returned false
at once, the \code{+} was copied as it was, and the lone blank stayed. The same
happened before \code{==}, \code{<} and \code{.and.} when their rules were off.

The close-delimiter handler had been relying on a later handler to tidy up after
it. The fix makes it stop leaving the blank before any operator character, whose
own rule decides the spacing:

\begin{lstlisting}[style=rust]
let skip_space = next_ch == ')' || next_ch == ']' || next_ch == ','
    || next_ch == '%' || next_ch == ':' || next_ch == '/' || next_ch == '*'
    || next_ch == ';' || (next_ch == '/' && self.peek(1) == Some(')'))
    // An operator is spaced by its own rule, which pops any blank
    // left here; with that rule off, a blank here would be lopsided
    || matches!(next_ch, '+' | '-' | '=' | '<' | '>' | '.');
if !skip_space {
    self.out.push(' ');
}
\end{lstlisting}

Two sibling rules in \code{needs\_space\_before\_delimiter} had the same flaw
from the other side, writing \code{arr= [1]} and \code{6]+ [1} at level 0:
the blank before an opening bracket after \code{=} or \code{+} is now written
only when the assignment or plus/minus rule is on. Output at the default level
did not change.

\begin{exercise}[Why each character?]
Remove the characters from the \code{matches!} above one at a time, and run
\code{cargo test} after each removal. Which test fails for each? For any
character whose removal no test notices, write a test that does.
\end{exercise}

\begin{exercise}[A new rule]
Add a twelfth whitespace rule, \code{power}, that spaces \code{**} at level 4
only. Chapter~\ref{ch:recipes} lists every place a new rule has to be
registered; follow it, and check that \code{--whitespace-power} appears in
\code{--help} without touching \file{cli.rs}. Why does that work?
\end{exercise}

% ---------------------------------------------------------------------------
\chapter{Long Lines and Comments}\label{ch:lines}
% ---------------------------------------------------------------------------

\section{Splitting long lines}\label{sec:linesplit}

Line splitting runs at step 12 of \code{format\_pass}, after indentation, on the
physical lines of one logical line. Its entry point, in
\file{format/line\_split.rs}, is

\begin{lstlisting}[style=rust]
pub fn split_long_lines(lines: &[String], indents: &[usize], max_line_length: usize,
                        indent_size: usize) -> (Vec<String>, Vec<usize>, Vec<usize>)
\end{lstlisting}

It returns the new lines, their indents, and for each new line the index of
the input line it came from, the \emph{origin}, which the comment logic needs.

For each line longer than the limit:

\begin{enumerate}
  \item If it has an inline comment, the comment is detached first
    (\code{split\_inline\_comment}). That may be enough on its own; if the code is
    still too long, it is split, and the comment goes on a line of its own
    after the fragments.
  \item Otherwise the line is split with \code{split\_with\_alignment}. If it
    cannot be split usefully, it is left as it is, over-length.
\end{enumerate}

\subsection{Where to break}

\code{find\_split\_position} finds a break within a given width. It collects
the positions of blanks and commas in code (never inside strings) and
prefers, from the right, a blank that leaves at least 12 characters for the
next line; failing that, a position after a comma that leaves more than 4.
\code{auto\_split\_line} uses it repeatedly: the first fragment must fit in
the limit minus the indent minus 2 (for the \code{ \&}), and each later
fragment in the limit minus the continuation indent. It gives up if the limit is
under 40, if the line already starts with \code{\&}, or if the line has a
comment.

Two decisions keep the result sensible:

\begin{itemize}
  \item If the tail of a statement cannot be broken (a long string, a long
    \code{a\%b\%c} path), the breaks already found are kept and the tail is
    emitted over-length. Before 0.5.0, the whole split was discarded.
  \item If, even so, some fragment is still over the limit, the split is only
    made if it actually saves width: if the widest fragment plus the
    continuation indent is at least as wide as the original line, the
    original is kept. Otherwise \code{call \&} in front of a 200-character
    call would buy a line and save one character.
\end{itemize}

\subsection{Splitting where the next pass would align}

A freshly split line must not move when it is formatted again, so its
continuation fragments have to go exactly where the aligner would put a
hand-written continuation. \code{split\_with\_alignment} asks the aligner:

\begin{lstlisting}[style=rust]
let guess = indent + indent_size;
let probe = auto_split_line(line, indent, max_line_length, guess)?;
let follow = indent + aligned_continuation_offset(line, &probe[0], indent_size);
if follow == guess {
    return Some((probe, guess));
}
auto_split_line(line, indent, max_line_length, follow)
    .map_or(Some((probe, guess)), |split| Some((split, follow)))
\end{lstlisting}

It splits once with a guessed continuation column (one indent deeper), asks
\code{F90Aligner} where it would put a continuation of that first fragment,
and, if that differs, splits again with the aligner's column, since a different
column changes how much fits on each fragment. If the statement already had
continuation lines, their column is known and used directly
(\code{follow\_indent}).

\code{LINE\_SPLIT\_THRESHOLD} (1024) caps the whole mechanism: with
\code{line\_length = 0}, or any limit of 1024 or more, nothing is split.

\section{Writing lines and placing comments}

\code{write\_output\_line} in \file{process.rs} writes one output line: the
label if it is the first line of a labelled statement, the code, and the
comment if this output line carries one. Its rules for comments are the most
intricate in fprettier, because they must be stable across the two passes.

\paragraph{Which line gets the comment.} \code{compute\_comment\_indices}
finds, for each origin, its last output line, and that line carries the
origin's comment. It also records which origins were split into several lines.

\paragraph{Detaching.} A comment moves to a new line when the code was split, or
when the code, the gap and the comment together exceed the limit:

\begin{lstlisting}[style=rust]
let inline_gap = if pass_ctx.config.normalize_comment_spacing {
    pass_ctx.config.comment_spacing
} else {
    original_gap.max(1)
};
let should_detach_comment = has_comment
    && !write_ctx.is_preprocessor
    && !is_comment_only
    && !comment_trimmed.is_empty()
    && write_ctx.effective_line_length < LINE_SPLIT_THRESHOLD
    && (write_ctx.split_origins.contains(&origin)
        || line_to_write.trim_end().len() + inline_gap + comment_trimmed.len()
            > write_ctx.effective_line_length);
\end{lstlisting}

A comment on a preprocessor line is never detached, just as line splitting
never breaks one. A fypp comment starts with \code{\#!}, and since
\code{CharFilter} reads the \code{!} as the start of a Fortran comment,
detaching it used to leave a bare \code{\#} on the directive and move the rest
of the comment to a line of its own, which fypp rejects.

The detached comment is indented at the line's computed indent.

\paragraph{The gap.} An inline comment that stays keeps \code{original\_gap},
the blanks at the end of the code part in \code{FortranLine::lines}
(Chapter~\ref{ch:reading}), unless spacing is normalised.

\paragraph{FORD comments.} A comment-only line beginning with \code{!!} is
written at its original indentation.

\paragraph{The one-space marker.} The subtlest piece is a one-character
string:

\begin{lstlisting}[style=rust]
// If the line is now empty (comment-only) but was originally indented,
// preserve one space as a marker so that Pass 2 of two-pass formatting
// can detect the original indentation. Only when an indent pass is
// actually coming: with indentation off nothing consumes the marker
// and the comment gains a space on every run.
if trimmed.is_empty()
    && pass_ctx.config.impose_indent
    && origin < write_ctx.lines_were_indented.len()
    && write_ctx.lines_were_indented[origin]
{
    " "
} else {
    trimmed
}
\end{lstlisting}

Pass 2 decides whether to indent a comment-only line by whether it was
indented in \emph{its} input, which is pass 1's output. Pass 1 must therefore
leave some trace of the original indentation; it leaves one blank. The
condition on \code{impose\_indent} is what keeps the result a fixed point: if
no indentation pass follows, nothing removes the marker, and every run would add
a blank.

\begin{keyidea}[The marker, once more]
A detached comment takes the computed indent of its line. In pass 1, which
computes no indents, that is the column the code started in, so a comment
detached from an \emph{unindented} line started at column 0. Pass 2 then saw an
unindented comment-only line, took it for one the author had put at column 0,
and left it there:
\begin{lstlisting}[style=before]
subroutine s
    x = 1
! a fairly long trailing comment that pushes the line past the limit of sixty
end subroutine
\end{lstlisting}
The fix uses the same marker: when pass 1 detaches a comment at column 0 and an
indentation pass is still to come, it writes one blank first, so pass 2 indents
the comment with its code, at column 4. This kind of interaction between passes
is what the idempotence tests of the next chapter are there to catch.
\end{keyidea}

\begin{exercise}[Splitting limits]
With \code{-l 50}, format a \code{call} with a single 60-character string
argument. Is the line split? Now add three short arguments before the string.
Explain both results using the ``saves width'' rule.
\end{exercise}

\begin{exercise}[Preprocessor comments]
Remove \code{!write\_ctx.is\_preprocessor} from \code{should\_detach\_comment},
build, and format \file{examples/in/test\_fypp.f90} with \code{-l 60}. Find
the damaged directives. Then explain why the comment of a C preprocessor line
such as \code{\#define X 1 ! note} must not be detached either.
\end{exercise}
% ---------------------------------------------------------------------------
\chapter{Optional Rewrites: Case, Operators, \texttt{use}}\label{ch:rewrites}
% ---------------------------------------------------------------------------

The three features of this chapter are off by default and independent of the
main passes. Each is a self-contained module with a single entry point.

\section{Case conversion}

\subsection{Settings}

\file{format/case\_convert.rs} defines the mode and the five categories:

\begin{lstlisting}[style=rust]
pub const CASE_KEYS: [&str; 5] = ["keywords", "procedures", "operators", "constants", "types"];

pub enum CaseMode {
    #[default]
    NoChange = 0,
    Lower = 1,
    Upper = 2,
}
\end{lstlisting}

\code{CaseMode} has a hand-written \code{Deserialize} implementation that
accepts exactly 0, 1 and 2, so a configuration file with \code{keywords = 3}
fails at the parse boundary. \code{CASE\_KEYS} fixes the order of
\flag{case}'s values, the directive's values and the \code{[case\_dict]} keys
in one place. \code{CASE\_FLAG\_VALUES} is \code{CASE\_KEYS.len() - 1}: the
\flag{case} flag takes only four values, because clap would take a fifth
value as a file name; \code{types} has its own flag.

\subsection{Token tables}

Five case-insensitive, word-bounded alternations of names classify tokens:
135 keywords, 207 intrinsic procedures, the five intrinsic modules, the
dotted operators, and 65 named constants from \code{iso\_fortran\_env} and
\code{iso\_c\_binding}. The lists are commented by the standard that
introduced each name, up to Fortran 2023. A test, \code{test\_tables\_reach\_f2023},
checks that names from each revision are present.

\subsection{The algorithm}

\code{convert\_case} splits the line into masked regions with
\code{CharFilter::code\_and\_comments} (so that a \code{!} inside a string is
not mistaken for a comment), leaves string parts alone, and splits each code part
into tokens: runs of alphanumerics, underscores and dots, and single other
characters. A post-processing step, \code{rejoin\_signed\_exponents}, glues
\code{1.0e}, \code{-}, \code{3} back into one token; without it, the signed
branch of the numeric patterns was unreachable and \code{1.0e-3} was not
converted while \code{1.0e3} was, a bug fixed in 0.4.0.

Each token is then classified by \code{convert\_token\_with\_context}, in this
order:

\begin{enumerate}
  \item a numeric literal: its exponent letter and kind suffix get the
    \code{types} mode (\code{2.e3\_real64} $\to$ \code{2.E3\_REAL64});
  \item a dotted operator or logical literal: \code{operators};
  \item a keyword: \code{keywords};
  \item an intrinsic module name: \code{procedures};
  \item an intrinsic procedure, but only if the next non-blank token is
    \code{(}: \code{procedures};
  \item a named constant: \code{constants};
  \item anything else: unchanged.
\end{enumerate}

A small piece of context tracking in \code{convert\_case} handles
\code{kind=\textit{name}}: the token after \code{kind =} gets the \code{types}
mode. That is why \code{real(kind=dp)} becomes \code{real(kind=DP)} but
\code{real(dp)} does not change.

\begin{pitfall}[Keywords are words, not statements]
The keyword test is purely lexical. A variable called \code{count} is not
converted (\code{count} is an intrinsic, and it is not followed by
\code{(}), but a variable called \code{result} or \code{value} is a
keyword to the table and is converted. Keep that in mind before enabling
\code{keywords} on a codebase with such names.
\end{pitfall}

\section{Relational operator replacement}

\code{replace\_relational\_operators} in \file{format/replacements.rs} is a
short, careful scanner. It first marks which byte positions are code, using
\code{CharFilter::code}, and then walks the line, trying the six pairs at each
code position:

\begin{lstlisting}[style=rust]
const OPERATORS: [(&str, &str); 6] = [
    (".le.", "<="),
    (".ge.", ">="),
    (".eq.", "=="),
    (".ne.", "/="),
    (".lt.", "<"),
    (".gt.", ">"),
];
\end{lstlisting}

The order puts \code{<=} before \code{<}; otherwise \code{a <= b} would become
\code{a .lt.= b}. When converting to Fortran style it skips \code{=>}, which
would otherwise become \code{=.gt.}.

\section{Sorting \texttt{use} statements}

\file{format/sort\_use.rs} is the only component that works on the raw text
of the whole file rather than on \code{FortranLine}s. It runs before both passes.

\subsection{Units and groups}

\code{split\_into\_units} groups physical lines into \code{Unit}s, one per
statement, following continuations (and comment-only or preprocessor lines inside
them):

\begin{lstlisting}[style=rust]
struct Unit {
    lines: Vec<String>,
    is_use: bool,
    /// Never moved and never rewritten (deactivation markers, preprocessor
    /// directives inside the statement, semicolons, unparseable module name).
    frozen: bool,
    /// Lowercased module name, empty unless `is_use`.
    key: String,
}
\end{lstlisting}

A unit is \emph{frozen} if it carries a \code{!\&} marker or is inside a
\code{!\&<}\ldots\code{!\&>} block, contains a preprocessor line, contains a
semicolon (a second statement), or is a \code{use} whose module name cannot be
parsed. \code{sort\_groups} then stable-sorts each maximal run of sortable units
by key. Anything that is not a sortable \code{use}, including a blank line,
ends a group, so \code{use} statements are never moved across a comment, a
blank line or a preprocessor branch.

Because units are whole physical lines, reordering is just moving line blocks:
indentation, line endings and trailing comments travel with their
statement.

\subsection{Sorting \texttt{only:} lists}

\code{sort\_only\_list} joins a unit's code into one line, finds the
\code{only:} list by walking the grammar (\code{only} is not a reserved word,
so a module called \code{only} must not confuse it), splits the list on
top-level commas (\code{operator(+)} contains none), sorts by the local name of
each item, and writes the statement back as \emph{one} line. If that line is
too long, the normal line splitter breaks it later. A unit with a comment on any
line but its last is left alone, because the comment would have nowhere to go.

\begin{exercise}[Keyword collisions]
Write a variable declaration and assignment using a variable called
\code{result}, and format it with \code{--case 2 0 0 0}. Then propose a rule
that would avoid converting it, and say what it would cost.
\end{exercise}

\begin{exercise}[Frozen units]
Construct a group of four \code{use} statements in which the third is frozen.
How many sorted groups does \code{sort\_groups} see? Check your answer with a
unit test.
\end{exercise}

% ---------------------------------------------------------------------------
\chapter{The Driver: \texttt{main.rs} and the CLI}\label{ch:driver}
% ---------------------------------------------------------------------------

\section{Arguments}

\file{cli.rs} builds the \code{clap} command with the builder API and converts
the matches into a plain struct, \code{CliArgs}, so that the rest of the
program never touches clap. Most of it is one \code{.arg(...)} call per flag.
The eleven whitespace flags are generated from the same table that defines the
rules:

\begin{lstlisting}[style=rust]
for (key, _, help) in WS_FLAGS {
    let name = ws_flag_name(key);   // "whitespace-" + key
    let mut arg = Arg::new(name.clone())
        .long(name)
        .help(help)
        .value_name("BOOL")
        .num_args(0..=1)
        .require_equals(true)
        .default_missing_value("true")
        .value_parser(clap::value_parser!(bool));
    // This flag was --whitespace-assignment before it was made to match
    // its config key; the old spelling still works, undocumented.
    if key == "assignments" {
        arg = arg.alias("whitespace-assignment");
    }
    cmd = cmd.arg(arg);
}
\end{lstlisting}

\code{WS\_FLAGS} in \file{config.rs} pairs each rule's \code{whitespace\_dict}
key with its index in the flags array and its help text, so that a flag, its
configuration key and the behaviour it controls cannot drift apart.
\code{require\_equals(true)} is what lets \code{-{}-whitespace-comma} be followed
by a file name without clap taking the file name as the flag's value.

\section{Building a configuration}

\code{build\_config} in \file{main.rs} starts from an explicit
\flag{config} file, or from the discovered files
(\code{Config::from\_discovered\_files}), then applies
the flags as Chapter~\ref{ch:config} describes. The one-way booleans are a
pair of operators:

\begin{lstlisting}[style=rust]
// Every boolean flag is one-way: a --no-* flag can only clear a setting the
// config file turned on, and the rest can only turn one on.
config.impose_indent &= !args.no_indent;
config.impose_whitespace &= !args.no_whitespace;
config.indent_fypp &= !args.no_indent_fypp;
config.indent_mod &= !args.no_indent_mod;
config.strict_indent |= args.strict_indent;
config.format_decl |= args.format_decl;
config.sort_use |= args.sort_use;
// ...
\end{lstlisting}

Discovery and merging live in \file{config.rs}.
\code{discover\_config\_files} returns the home file and the ancestors' files,
least specific first. \code{from\_discovered\_files} merges their TOML tables
recursively with \code{merge\_tables} and deserialises after \emph{each} file,
not just at the end, so that an unknown key or a bad value is reported against
the file that introduced it. \code{Config} itself is deserialised with
\code{\#[serde(default, deny\_unknown\_fields)]}, which is what makes every key
optional and every unknown key an error.

\section{Processing files}

\code{main} decides between standard input and files, then:

\begin{enumerate}
  \item reports each input path that does not exist (and remembers to fail);
  \item collects files with \code{collect\_files} (\code{walkdir} with
    \code{follow\_links(true)}, a depth limit of 256 when recursive and 1
    otherwise, and the extension and exclusion filters);
  \item decides whether to run sequentially: for \flag{stdout}, \flag{diff},
    \flag{check} and \code{-j 1};
  \item if parallel, pre-warms the regular expressions;
  \item calls \code{process\_files}.
\end{enumerate}

\subsection{A configuration per directory}

Without \flag{config}, each file may have its own configuration, since
\file{fprettier.toml} files can sit at any level. Discovering it costs a
\code{stat} per ancestor directory, so \code{process\_files} builds one
configuration per distinct parent directory before starting the parallel work:

\begin{lstlisting}[style=rust]
let dir_configs: HashMap<PathBuf, Config> = if base_config.is_none() {
    let mut dirs: Vec<PathBuf> = files.iter()
        .filter_map(|f| f.parent().map(Path::to_path_buf)).collect();
    dirs.sort();
    dirs.dedup();
    dirs.into_iter()
        .filter_map(|dir| build_config(args, Some(&dir)).ok().map(|c| (dir, c)))
        .collect()
} else {
    HashMap::new()
};
\end{lstlisting}

A directory whose configuration failed to build has no entry; when a file from
it is processed, \code{build\_config} is called again for that file, so the
error is reported against the file. Each file is then processed by a closure
run with \code{files.par\_iter().for\_each(run)} or, sequentially,
\code{files.iter().for\_each(run)}; two \code{AtomicUsize} counters
collect the number of changed files and errors.

\subsection{One file}

\code{process\_single\_file} checks the size (100\,MB) before reading, reads the
file, checks \flag{exclude-max-lines}, and looks for a directive with
\code{find\_directive}. Only when a directive is found is the configuration
cloned and overridden; most files have none, and cloning a \code{Config} with
its maps for every file of a large tree is measurable. It then calls
\code{format\_file}, compares the output with the input, and acts on the
output mode.

\subsection{Writing safely}

\code{write\_in\_place} implements the atomic write promised in
Chapter~\ref{ch:start}:

\begin{lstlisting}[style=rust]
let target = std::fs::canonicalize(path)?;
let temp = target.with_file_name(format!(
    ".{}.fprettier{}-{}",
    name.to_string_lossy(),
    std::process::id(),
    COUNTER.fetch_add(1, Ordering::Relaxed)
));
let write_and_rename = || -> std::io::Result<()> {
    std::fs::write(&temp, contents)?;
    // rename() does not carry the target's mode over, so copy it across
    std::fs::set_permissions(&temp, std::fs::metadata(&target)?.permissions())?;
    std::fs::rename(&temp, &target)
};
write_and_rename().inspect_err(|_| {
    let _ = std::fs::remove_file(&temp);
})?;
\end{lstlisting}

\code{canonicalize} resolves symbolic links, so the rename replaces the file
the link points to rather than the link. The process ID and a counter make the
temporary name unique across threads and concurrent runs.

\subsection{Writing to a closed pipe}

Rust ignores \code{SIGPIPE}, so writing to a pipe whose reader has exited
returns a \code{BrokenPipe} error, and \code{print!} panics on it. All of
fprettier's standard output goes through \code{write\_stdout}, which exits
quietly instead:

\begin{lstlisting}[style=rust]
fn write_stdout(bytes: &[u8]) -> Result<()> {
    match io::stdout().write_all(bytes) {
        Err(e) if e.kind() == io::ErrorKind::BrokenPipe => std::process::exit(0),
        result => Ok(result?),
    }
}
\end{lstlisting}

\section{Performance}

Three measured decisions are recorded in comments in \file{main.rs}, each with
its benchmark:

\begin{itemize}
  \item \textbf{mimalloc} as the global allocator is 5.5--9.5\% faster than the
    system allocator at every thread count on a 3000-file tree. The gain is
    largest single-threaded, so it comes from cheaper small allocations
    (formatting creates many short-lived \code{String}s), not from less lock
    contention.
  \item \textbf{Pre-warming} the regular expressions saves 10--40\,ms on such
    a run.
  \item \textbf{Per-directory configuration} avoids about ten \code{stat}
    calls per file.
\end{itemize}

Where the time goes inside \code{format\_file} has not been recorded in the
source; finding out is the next exercise.

\begin{exercise}[Measure]
Run \code{cargo bench}. The \code{cli/tree 480 files} benchmark copies
\file{examples/in} into a temporary tree and formats it. Then replace mimalloc
with the system allocator (delete the \code{\#[global\_allocator]} static), run
it again, and compare with the table in \file{main.rs}.
\end{exercise}

\begin{exercise}[Profile]
Profile \code{format\_file} on a large input (concatenate
\file{examples/in/*.f90} a few hundred times), for example with
\code{cargo flamegraph -{}-bench bench}. Which functions dominate? A good
guess beforehand: the indenter tries every enabled scope's patterns against
every statement. Was the guess right?
\end{exercise}

\begin{exercise}[Parallel output]
Why must \flag{diff} run sequentially, when \flag{check} only prints one line per
file? Make \flag{check} run in parallel, collecting the names and printing them
sorted at the end. Is the result still deterministic?
\end{exercise}
% ---------------------------------------------------------------------------
\chapter{Testing fprettier}\label{ch:testing}
% ---------------------------------------------------------------------------

\section{The four kinds of test}

\code{cargo test} runs 385 tests in four groups:

\begin{center}
\small
\begin{tabularx}{\linewidth}{@{}l r X@{}}
\toprule
Where & Tests & What they check\\
\midrule
\code{\#[cfg(test)]} modules in \file{src/} & 222 & One function at a time: a regex, a \code{CharFilter} state, an aligner column, a whitespace handler\\
\file{tests/integration\_test.rs} & 70 & Components together through the public API: the indenter fed statement by statement, \code{format\_file} on small programs\\
\file{tests/compatibility\_test.rs} & 90 & Whole files and focused snippets against expected output, including the \file{examples/in} $\to$ \file{examples/out} pairs\\
\file{tests/cli\_test.rs} & 3 & The binary itself: \flag{check} and \flag{diff} exit codes, output and the promise not to touch files\\
\bottomrule
\end{tabularx}
\end{center}

\subsection{Golden files}

\file{examples/in} and \file{examples/out} hold eight pairs of real-world-sized
files: a general example, fypp, \code{associate}, \code{where}/\code{forall},
\code{namelist}/\code{block}/\code{select}, invalid code, case conversion
and \code{use} sorting. Each has a test that formats the input with a specific
configuration and compares with the output:

\begin{lstlisting}[style=rust]
#[test]
fn test_example_compatibility() {
    let input = get_example_path("in", "example.f90");
    let expected = get_example_path("out", "example.f90");
    // Default settings - format_decl=false preserves existing :: alignment
    let config = Config {
        impose_whitespace: true,
        impose_indent: true,
        indent: 3,
        whitespace: 2,
        format_decl: false,
        ..Default::default()
    };
    test_compatibility(&input, &expected, &config);
}
\end{lstlisting}

\code{test\_compatibility} does not \code{assert\_eq!} two whole files, which
would print two walls of text; it reports the first line that differs,
then a line-count difference, and finally the case where only the trailing
newline differs.

\begin{pitfall}[Do not regenerate expected output blindly]
When a change makes a golden-file test fail, it is tempting to overwrite
\file{examples/out} with the new output. Do that only after reading every
changed line and convincing yourself that each one is an improvement. The
golden files encode years of decisions, many of them about compatibility with
fprettify; a diff in them is a behaviour change for every user.
\end{pitfall}

\subsection{Testing idempotence}

The fixed-point property of Chapter~\ref{ch:tour} is tested directly in places
where it once failed: format, format the result, and compare.

\begin{lstlisting}[style=rust]
let result1 = run_format(input, &config);
let result2 = run_format(&result1, &config);
assert_eq!(result1, result2, "OMP conditional continuation should be idempotent");
\end{lstlisting}

It is also implied by the golden files: every \file{examples/out} file is a
fixed point under its test's configuration. You can check this from the shell;
for example, \code{fprettier -{}-check -i 3 examples/out/example.f90} passes.

\section{Continuous integration}

\file{.github/workflows/ci.yml} runs three jobs on every push to \code{main}
and every pull request:

\begin{description}[font=\sffamily\small\bfseries]
  \item[Test] \code{cargo test -{}-all-features} on stable.
  \item[Lint] \code{cargo +nightly fmt -{}-check} (nightly, because
    \file{rustfmt.toml} uses \code{group\_imports} and
    \code{imports\_granularity}, which stable rustfmt ignores);
    \code{cargo +stable clippy -{}-all-targets -{}-all-features -{}- -D warnings},
    with \code{clippy::pedantic} enabled in \file{Cargo.toml}; and
    \code{cargo doc} with warnings as errors.
  \item[MSRV] \code{cargo check} on Rust 1.87.0, the \code{rust-version}
    declared in \file{Cargo.toml}. The comment there names the feature that sets
    the floor: \code{is\_multiple\_of} on unsigned integers, used by the
    indentation check in \code{inspect\_file}.
\end{description}

Run the same checks locally before you push:

\begin{lstlisting}[style=shell]
cargo test
cargo +nightly fmt --check
cargo clippy --all-targets --all-features -- -D warnings
RUSTDOCFLAGS="-D warnings" cargo doc --no-deps
\end{lstlisting}

\section{Benchmarks}

\file{benches/bench.rs} has two Criterion benchmarks: \code{format\_file} on
\file{examples/in/example.f90}, which measures the formatter alone, and the CLI on
a temporary tree of 480 files, which measures discovery, parallelism and I/O
too. The tree benchmark passes an empty \flag{config} so that a
\file{fprettier.toml} in your home directory cannot skew it.

\begin{exercise}[A test that should exist]
Until this book was written, no test failed if \code{inspect\_file} read the
wrong buffer (Exercise~9.3). Choose another invariant that this book has described (for
example: blank lines are transparent to stacking; \code{--check} never writes
files; a labelled \code{do} closes at its label when two loops share it) and
find out whether a test pins it down. If none does, write one.
\end{exercise}

\begin{exercise}[An idempotence sweep]
Write a test that formats every file in \file{examples/in} with the default
configuration, formats the result again, and asserts that the second pass
changes nothing. Does it pass? Try it with \code{-l 60} and with
\code{whitespace: 4} too.
\end{exercise}

% ---------------------------------------------------------------------------
\chapter{Recipes for Contributors}\label{ch:recipes}
% ---------------------------------------------------------------------------

This chapter collects step-by-step instructions for the changes contributors
most often make. Each lists every place that needs touching; the compiler will
catch some omissions, but not all.

\section{Fixing a formatting bug}

\begin{enumerate}
  \item \textbf{Reproduce} with \code{fprettier -s} and the reporter's
    configuration. Use \flag{config} with an empty file, plus explicit flags,
    so your own \file{fprettier.toml} does not interfere.
  \item \textbf{Minimise.} Delete lines until the bug disappears, then put the
    last one back. Most bugs reproduce in three to five lines.
  \item \textbf{Locate.} Run with \flag{no-indent} and then with
    \flag{no-whitespace}, to learn which pass is responsible. If the bug
    involves continuations, check whether the lines have a leading \code{\&}
    (manual alignment, Section~\ref{sec:manual}).
  \item \textbf{Write the failing test first}, in
    \file{tests/compatibility\_test.rs} if it is about output, or as a unit
    test next to the function you suspect.
  \item \textbf{Fix} the narrowest piece of code that explains the bug, and
    leave a comment saying \emph{why} the code is the way it is. The existing
    comments in fprettier are mostly of this kind; follow their example.
  \item \textbf{Check idempotence}: format the test's output again.
  \item \textbf{Run the full suite} and the lint checks of
    Chapter~\ref{ch:testing}, and read any change in \file{examples/out}.
  \item \textbf{Record it} under \emph{Unreleased} in \file{CHANGELOG.md},
    describing the behaviour users saw, not the code you changed.
\end{enumerate}

\section{Adding a whitespace rule}

\begin{enumerate}
  \item In \file{config.rs}, add a \code{WS\_*} index constant and a row in
    \code{WS\_FLAGS} (key, index, help text). The \flag{whitespace-\textit{key}}
    flag and the \code{[whitespace\_dict]} key come from that row
    automatically.
  \item Grow the flags array from \code{[bool; 11]} to \code{[bool; 12]}.
    The type appears in \file{config.rs}, \file{process.rs} and
    \file{format/whitespace.rs}; the compiler finds every occurrence.
  \item Add the rule's value to each level preset in
    \code{get\_whitespace\_flags}.
  \item Implement the rule: a new \code{try\_*} handler in \code{Charwise}
    (mind its position in the chain), or a step in
    \code{add\_whitespace\_context}.
  \item Test it: a unit test in \file{whitespace.rs}, and check the flags-array
    tests in \file{config.rs}, which list each level's array literally.
  \item Document it in the README's \code{[whitespace\_dict]} block and in
    \file{CHANGELOG.md}.
\end{enumerate}

\section{Adding a scope}

\begin{enumerate}
  \item In \file{parser/patterns.rs}, add the opener and closer patterns, and a
    continuation pattern if the construct has one. Follow the existing shape:
    \code{SOL\_STR}, an optional construct name, \code{\textbackslash s*}
    between keywords that may be run together, \code{EOL\_STR}.
  \item Append a row to \code{SCOPES} in \file{scope.rs}. Never insert in the
    middle: the indices are the scopes' identity.
  \item Decide whether the construct belongs in \code{opens\_indented\_scope}
    in \file{process.rs}, that is, whether an \code{if} written flush with it
    should be indented rather than stacked. For a block construct the answer is
    usually yes.
  \item If its \code{end} can be written run together (\code{endfoo}), add the
    keyword to \code{END\_RE} and the two \code{END\_*} patterns in
    \file{whitespace.rs}, and to the keyword table in \file{case\_convert.rs}.
  \item Test the patterns in \file{patterns.rs}, the indentation in
    \file{tests/integration\_test.rs} (the \code{indents\_for} helper makes
    this one line), and a whole construct in \file{compatibility\_test.rs}.
\end{enumerate}

\section{Adding a configuration option}

\begin{enumerate}
  \item Add the field to \code{Config} and its default to
    \code{impl Default for Config}. Because of \code{\#[serde(default)]}, the
    TOML key exists from then on.
  \item Add a field to \code{CliArgs}, an \code{Arg} in \code{build\_cli}, and
    the conversion in \code{args\_from\_matches}.
  \item Apply it in \code{build\_config} in \file{main.rs}, respecting the
    one-way convention for booleans.
  \item If directives should accept it, extend \code{DirectiveOverrides},
    \code{parse\_directive\_args} and \code{apply\_directive\_overrides}.
  \item If it has invalid values, check them in \code{Config::validate}.
  \item Use it, document it in the README's key list, and add it to the CHANGELOG.
\end{enumerate}

\section{Style}

The code base follows a few conventions worth keeping:

\begin{itemize}
  \item \code{clippy::pedantic} is on; allow a lint locally with a reason
    rather than globally.
  \item Comments explain why, often with the bug that made the code necessary.
    When you remove a special case, find out first which test it serves.
  \item References to the Fortran standard use clause numbers of Fortran 2023
    (\emph{6.3.2.1}) or rule numbers (\emph{R1002}), so they can be checked.
  \item Numbers that drove a decision (benchmark results, limits) are written
    next to the code they justify, with the date they were measured.
\end{itemize}

\begin{exercise}[Your first contribution]
At short line limits, fprettier's output is not always a fixed point. Formatting
\file{examples/in/example.f90} with \code{-l 60} and then formatting the result
again moves the pieces of a statement that was split inside nested brackets:
\begin{lstlisting}[style=plain]
first run:                                      r, k, l) + i
second run:                                               r, k, &
                                                          l) + i
\end{lstlisting}
Four of the eight example files do this at \code{-l 60} or \code{-l 40}. Follow
the bug-fixing recipe from start to finish: write the idempotence sweep of
Exercise~16.2 as a failing test, find out why the column chosen by
\code{split\_with\_alignment} differs from the one the aligner computes on the
next run, fix it, and add the CHANGELOG entry.
\end{exercise}

% ===========================================================================
\appendix
% ===========================================================================

\chapter{Configuration Reference}\label{app:config}

\begin{center}
\small
\begin{tabularx}{\linewidth}{@{}l l l X@{}}
\toprule
Key & Type & Default & Meaning\\
\midrule
\code{indent} & integer & 4 & Spaces per indent level; at least 1\\
\code{line\_length} & integer & 132 & Maximum line length; at most 10\,000; 0 for none\\
\code{whitespace} & 0--4 & 2 & Whitespace level (Table~\ref{tab:ws-levels})\\
\code{impose\_indent} & bool & true & Run the indentation pass\\
\code{impose\_whitespace} & bool & true & Run the whitespace pass\\
\code{strict\_indent} & bool & false & Disable the stacking rule\\
\code{indent\_fypp} & bool & true & fypp blocks open scopes and are indented\\
\code{indent\_mod} & bool & true & Indent program, module and submodule bodies\\
\code{format\_decl} & bool & false & Normalise spacing around \code{::}\\
\code{comment\_spacing} & integer & 1 & Gap before inline comments when normalised\\
\code{normalize\_comment\_spacing} & bool & false & Apply \code{comment\_spacing} to every inline comment\\
\code{enable\_replacements} & bool & false & Rewrite relational operators\\
\code{c\_relations} & bool & false & With replacements: C style, not Fortran style\\
\code{sort\_use} & bool & false & Sort \code{use} statements within a group\\
\code{sort\_use\_only} & bool & false & Sort \code{only:} lists\\
\code{[whitespace\_dict]} & table & empty & Per-rule overrides: \code{comma}, \code{assignments}, \code{relational}, \code{logical}, \code{plusminus}, \code{multdiv}, \code{print}, \code{type}, \code{intrinsics}, \code{decl}, \code{concat}\\
\code{[case\_dict]} & table & empty & 0/1/2 for \code{keywords}, \code{procedures}, \code{operators}, \code{constants}, \code{types}\\
\bottomrule
\end{tabularx}
\end{center}

\chapter{Command-Line Reference}\label{app:cli}

\begin{center}
\small
\begin{tabularx}{\linewidth}{@{}l X@{}}
\toprule
Flag & Effect\\
\midrule
\code{-i}, \flag{indent} \textit{N} & Spaces per indent level\\
\code{-l}, \flag{line-length} \textit{N} & Maximum line length\\
\code{-w}, \flag{whitespace} \textit{N} & Whitespace level 0--4\\
\flag{whitespace-\textit{rule}}[\code{=}\textit{BOOL}] & Override one whitespace rule\\
\flag{no-indent}, \flag{no-whitespace} & Skip a main pass\\
\flag{strict-indent} & Disable the stacking rule\\
\flag{no-indent-fypp}, \flag{no-indent-mod} & See \code{indent\_fypp}, \code{indent\_mod}\\
\flag{format-decl} & Normalise \code{::} spacing\\
\flag{normalize-comment-spacing}, \flag{comment-spacing} \textit{N} & Inline comment gap\\
\flag{case} \textit{K P O C} & Case of keywords, procedures, operators, constants (0/1/2)\\
\flag{case-types} \textit{N} & Case of kind values, kind suffixes, exponent letters\\
\flag{enable-replacements}, \flag{c-relations} & Relational operator spelling\\
\flag{sort-use}, \flag{sort-use-only} & Sort \code{use} statements, \code{only:} lists\\
\code{-s}, \flag{stdout} & Print instead of writing files\\
\code{-d}, \flag{diff} & Print a unified diff\\
\flag{check} & List files that would change; exit 1 if any\\
\code{-c}, \flag{config} \textit{FILE} & Use this configuration file only\\
\code{-r}, \flag{recursive} & Recurse into directories\\
\code{-e}, \flag{exclude} \textit{GLOB} & Exclude matching paths (repeatable)\\
\code{-f}, \flag{fortran} \textit{EXT} & Extra Fortran extension (repeatable)\\
\code{-m}, \flag{exclude-max-lines} \textit{N} & Skip files longer than \textit{N} lines\\
\code{-j}, \flag{jobs} \textit{N} & Worker threads (1: sequential)\\
\code{-D}, \flag{debug} & Print configuration diagnostics\\
\code{-S}, \flag{silent} & No summary or skip messages\\
\bottomrule
\end{tabularx}
\end{center}

Inputs are files or directories; none, or \code{-}, reads standard input.

\chapter{The Scope Table}\label{app:scopes}

\begin{center}
\small
\footnotesize
\begin{tabular}{@{}r l l l l@{}}
\toprule
\# & Group & Opens & Continues & Closes\\
\midrule
0 & Core & \code{if (\ldots) then} & \code{else}, \code{else if} & \code{end if}\\
1 & Core & \code{do}, \code{do while}, \code{do concurrent} & & \code{end do}\\
2 & Core & \code{select case/rank/type} & \code{case}, \code{rank}, \code{type is}, \ldots & \code{end select}\\
3 & Core & \code{subroutine} & \code{contains} & \code{end subroutine}\\
4 & Core & \code{function} & \code{contains} & \code{end function}\\
5 & Core & \code{interface} & & \code{end interface}\\
6 & Core & \code{type} definition & \code{contains} & \code{end type}\\
7 & Core & \code{enum} & & \code{end enum}\\
8 & Core & \code{associate} & & \code{end associate}\\
9 & Core & (none) & & bare \code{end}, closes any\\
10 & Core & \code{block}, \code{block data} & & \code{end block}\\
11 & Core & \code{where (\ldots)} block & \code{elsewhere} & \code{end where}\\
12 & Core & \code{forall (\ldots)} block & & \code{end forall}\\
13 & Module & \code{module} & \code{contains} & \code{end module}\\
14 & Module & \code{submodule} & \code{contains} & \code{end submodule}\\
15 & Module & \code{program} & \code{contains} & \code{end program}\\
16 & Fypp & \code{\#:def} & & \code{\#:enddef}\\
17 & Fypp & \code{\#:if} & \code{\#:elif}, \code{\#:else} & \code{\#:endif}\\
18 & Fypp & \code{\#:for} & & \code{\#:endfor}\\
19 & Fypp & \code{\#:block} & & \code{\#:endblock}\\
20 & Fypp & \code{\#:call} & & \code{\#:endcall}\\
21 & Fypp & \code{\#:mute} & & \code{\#:endmute}\\
22 & Core & \code{critical} & & \code{end critical}\\
23 & Core & \code{change team} & & \code{end team}\\
24 & Core & \code{enumeration type} & & \code{end enumeration type}\\
\bottomrule
\end{tabular}
\end{center}

\chapter{Hints for Selected Exercises}

\paragraph{1.2} A missing path is an error (exit 1). An empty directory is not:
fprettier prints \code{No Fortran files found to format.} and exits 0.

\paragraph{2.1} \code{print *, 'a'//'b'}, \code{do while (i < 10)},
\code{x(1:n) = y(1:n)*2}, \code{p => t}. The concatenation stays tight below
level 4, and the colon of an array section is never spaced.

\paragraph{2.2} Only \code{SIZE(arr)} changes. The declared variable and the
assignment target are not followed by \code{(}, so they are not procedure
calls.

\paragraph{3.2} Nothing is sorted: directives do not accept \flag{sort-use}, and
unknown directive options are ignored.

\paragraph{4.2} At level 3 the output is \code{y = a * -b + c(-1) - 2e-3 - x\%e - 3}.
The \code{-b} and \code{(-1)} signs follow an operator and an open bracket, so
they are unary; \code{2e-3} is an exponent; and \code{x\%e-3} is not, because
\code{e} is a component name.

\paragraph{5.1} Indented by four, the loops are kept stacked. Indented by
three, column 3 is not a multiple of 4, so stacking is off and each loop is
indented.

\paragraph{5.2} No. Blank lines are transparent to \code{inspect\_file}
(Section~\ref{sec:inspect}).

\paragraph{6.1} (a) \code{!\&<} \ldots{} \code{!\&>}, or a leading \code{\&}
on each continuation line; (b) \flag{exclude}; (c) a trailing \code{!\&}; (d)
\flag{exclude-max-lines}.

\paragraph{10.2} \code{end associate} is written at the loop body's level
and changes nothing: \code{ENDASSOCIATE\_RE} matches, the popped scope is the
\code{do}, not an \code{associate}, so the \code{else} branch pushes it back
and the \code{end} is marked invalid, which \code{calculate\_indent} treats
like an ordinary statement. Without the push-back, the \code{do} would be gone
from \code{scope\_storage} while its indent stayed on \code{indent\_storage}, and
the stacks would be out of step for the rest of the file.

\paragraph{11.1} The \code{(} of \code{build} is at relative column 27, so the
continuation starts at $4 + 28 = 32$. The pushes are: the \code{=>} (column 7),
\code{make\_node(} (17), \code{build(} (28).

\paragraph{13.1} The single-string call is not split: breaking after
\code{call} would save almost nothing, and the string cannot be broken. With
three arguments before it, breaking after \code{c,} saves enough to be worth
it, even though the string line is still over 50 characters.

\backmatter
\chapter*{Afterword}
\addcontentsline{toc}{chapter}{Afterword}

A formatter is judged by what it leaves alone as much as by what it changes. The
reader who has come this far has seen how much of fprettier exists for that
purpose: a character filter so that strings and comments are never touched,
a pre-pass that reads how the author laid out loops, a single blank kept as a
marker so that a comment's indentation survives two passes, a split that asks
the aligner where the next run would put it. None of these is complicated on
its own. Together they make a tool that can run over a million lines of
someone else's Fortran and leave every one of them better, or exactly as it
was.

Contributions are welcome at \url{https://github.com/fredjones42/fprettier}.

\end{document}
